% Section file for Chapter: Twisted Holography and Quantum Gravity
% Result (G2): Gravitational Complexity
%
% The shadow depth r_max(A), a purely algebraic invariant of the
% boundary chiral algebra, determines the minimal homotopy-algebraic
% complexity of the bulk gravitational theory.  The four depth classes
% G/L/C/M emerge as the only possibilities from the structure of the
% obstruction tower.

\section{Gravitational complexity}
% label removed: sec:thqg-gravitational-complexity
\index{gravitational complexity|textbf}
\index{shadow depth!gravitational interpretation}

% ======================================================================
%  §1.  THE OBSTRUCTION TOWER IN THE GRAVITATIONAL SETTING
% ======================================================================

\subsection{The obstruction tower}
% label removed: subsec:thqg-obstruction-tower

The shadow obstruction tower
$\Theta_\cA^{\leq 2} \to \Theta_\cA^{\leq 3} \to
\Theta_\cA^{\leq 4} \to \cdots$
is the sequence of arity projections of the bar-intrinsic MC element
$\Theta_\cA = D_\cA - \dzero \in \MC(\gAmod)$
(Theorem~\ref*{V1-thm:mc2-bar-intrinsic}).  The MC equation at
arity~$r{+}1$ reads
\begin{equation}% label removed: eq:thqg-obstruction-recursion
d_2\,\Theta^{\leq r+1}_{r+1}
\;=\;
-\bigl[\Theta_\cA^{\leq r},\, \Theta_\cA^{\leq r}\bigr]_{r+1}
\;-\;
\hbar\,\Lambda_P\bigl(\Theta_\cA^{\leq r}\bigr)_{r+1},
\end{equation}
where $d_2 := d + [\Theta^{\leq 2}, -]$ is the twisted differential
and $\Lambda_P\colon \operatorname{Sym}^r(V_\cA^*) \to
\operatorname{Sym}^{r-2}(V_\cA^*)$ is the genus-loop operator
(Definition~\ref{def:nms-genus-loop-operator}).

\begin{definition}[Gravitational obstruction class]
% label removed: def:thqg-gravitational-obstruction
\index{gravitational obstruction class|textbf}
\index{obstruction class!gravitational}
The \emph{gravitational obstruction class at arity~$r{+}1$} is
\begin{equation}\label{eq:thqg-obstruction-class}
\mathfrak{o}^{(r+1)}_{\mathrm{grav}}(\cA)
\;:=\;
\bigl[\Theta_\cA^{\leq r},\, \Theta_\cA^{\leq r}\bigr]_{r+1}
\;+\;
\hbar\,\Lambda_P\bigl(\Theta_\cA^{\leq r}\bigr)_{r+1}
\;\in\;
Z^2\bigl(F^{r+1}\gAmod / F^{r+2}\gAmod,\, d_2\bigr).
\end{equation}
This is a cocycle: $d_2\,\mathfrak{o}^{(r+1)} = 0$ by the graded
Jacobi identity applied to the weight-$(r{+}2)$ component of the
MC equation.  The cohomology class
\begin{equation}% label removed: eq:thqg-obstruction-cohomology
o_{r+1}(\cA)
\;:=\;
[\mathfrak{o}^{(r+1)}_{\mathrm{grav}}(\cA)]
\;\in\;
H^2\bigl(F^{r+1}\gAmod / F^{r+2}\gAmod,\, d_2\bigr)
\end{equation}
vanishes if and only if the tower extends from arity~$r$ to
arity~$r{+}1$ without new data.
\end{definition}

\begin{proposition}[Two sources of gravitational obstruction]
% label removed: prop:thqg-two-sources
\ClaimStatusProvedHere
\index{gravitational obstruction!bracket and loop decomposition}
The obstruction~\eqref{eq:thqg-obstruction-class} decomposes into
two structurally independent terms:
\begin{enumerate}[label=\textup{(\roman*)}]
\item The \emph{bracket term}
  $\bigl[\Theta_\cA^{\leq r},\,\Theta_\cA^{\leq r}\bigr]_{r+1}$:
  a sum over connected stable graphs of total arity~$r{+}1$ with
  at least two vertices, each labelled by a transferred operation
  $\ell^{\mathrm{tr}}_{|v|}$ from the cyclic minimal model, with
  edges contracted by the propagator $P_\cA = H_\cA^{-1}$.
  This is the genus-$0$ contribution.
\item The \emph{genus-loop term}
  $\hbar\,\Lambda_P(\Theta_\cA^{\leq r})_{r+1}$:
  the BV operator $\Lambda_P$ applied to the lower-arity shadow,
  converting two external legs into an internal loop.  This increases
  genus by one.
\end{enumerate}
The bracket term is controlled by $L_\infty$-formality of
$\gAmod$; the genus-loop term by the cyclic structure.
\end{proposition}

\begin{proof}
Connected graphs contributing to arity~$r{+}1$ are of two types:
multi-vertex graphs (all edges internal, total arity = sum of
vertex arities minus number of internal edges), and single-vertex
graphs with one self-loop.  These are the only contributions to
the MC equation at weight~$r{+}1$ in the modular convolution
algebra (Lemma~\ref{lem:graph-sum-truncation}).
\end{proof}

\begin{remark}[The first four obstruction classes]
% label removed: rem:thqg-first-four
\index{obstruction class!explicit low-arity}
Expanding~\eqref{eq:thqg-obstruction-class}:
\begin{align}
o_3(\cA)
  &= [\mathrm{Sh}_2, \mathrm{Sh}_2]_3,
  % label removed: eq:thqg-o3 \\
o_4(\cA)
  &= [\mathrm{Sh}_2, \mathrm{Sh}_3]_4
  + \tfrac{1}{2}[\mathrm{Sh}_3, \mathrm{Sh}_3]_4,
  % label removed: eq:thqg-o4 \\
o_5(\cA)
  &= [\mathrm{Sh}_2, \mathrm{Sh}_4]_5
  + [\mathrm{Sh}_3, \mathrm{Sh}_3]_5
  + [\mathrm{Sh}_3, \mathrm{Sh}_4]_5
  + [\mathrm{Sh}_4, \mathrm{Sh}_3]_5,
  % label removed: eq:thqg-o5 \\
o_6(\cA)
  &= [\mathrm{Sh}_2, \mathrm{Sh}_5]_6
  + [\mathrm{Sh}_3, \mathrm{Sh}_4]_6
  + [\mathrm{Sh}_3, \mathrm{Sh}_5]_6
  + \tfrac{1}{2}[\mathrm{Sh}_4, \mathrm{Sh}_4]_6
  + [\mathrm{Sh}_5, \mathrm{Sh}_3]_6.
  % label removed: eq:thqg-o6
\end{align}
The pattern: $o_{r+1}$ is a polynomial in the lower shadows,
with the bracket encoding pairwise sewing through the propagator.
\end{remark}


% ======================================================================
%  §2.  SHADOW DEPTH AND THE COMPLEXITY INVARIANT
% ======================================================================

\subsection{Shadow depth as a gravitational invariant}
% label removed: subsec:thqg-shadow-depth-grav

\begin{definition}[Gravitational complexity]
% label removed: def:thqg-gravitational-complexity
\index{gravitational complexity!definition|textbf}
\index{shadow depth!as gravitational invariant}
Let $\cA$ be a cyclically admissible chiral algebra
(Definition~\ref{def:cyclically-admissible}).
The \emph{gravitational complexity} of~$\cA$ is
\begin{equation}% label removed: eq:thqg-grav-complexity
\rho_{\mathrm{grav}}(\cA)
\;:=\;
r_{\max}(\cA)
\;=\;
\sup\{r \geq 2 : \cA^{\mathrm{sh}}_{r,0} \neq 0\}
\;\in\;
\{2, 3, 4, \ldots\} \cup \{\infty\},
\end{equation}
where $\cA^{\mathrm{sh}}_{r,0} =
H^*(\Defcyc(\cA))_{r,0}$ is the arity-$r$, genus-$0$
component of the shadow algebra
(Definition~\ref{def:shadow-algebra}).

Two related homotopy-theoretic invariants are the $A_\infty$-depth
$d_\infty(\cA)$ and the $L_\infty$-formality level $f_\infty(\cA)$.
Equivalently, by the operadic complexity theorem of Vol~I
\textup{(}Theorem~\textup{\ref*{V1-thm:operadic-complexity-detailed}}\textup{)}:
\begin{equation}% label removed: eq:thqg-three-invariants
\rho_{\mathrm{grav}}(\cA)
\;=\;
d_\infty(\cA)
\;=\;
f_\infty(\cA),
\end{equation}
where $d_\infty$ is the $A_\infty$-depth and $f_\infty$
the $L_\infty$-formality level.
\end{definition}

\begin{theorem}[Gravitational complexity controls bulk vertices]
% label removed: thm:thqg-complexity-controls-bulk
\ClaimStatusProvedHere
\index{gravitational complexity!bulk vertex theorem}
Let $\cA$ be a chirally Koszul logarithmic $\SCchtop$-algebra
\textup{(}Definition~\textup{\ref{def:log-SC-algebra}}\textup{)}.
The shadow obstruction tower stabilises at arity $r_{\max}(\cA)$:
\begin{equation}% label removed: eq:thqg-stabilization
\Theta_\cA^{\leq r}
\;=\;
\Theta_\cA^{\leq r_{\max}}
\quad \text{for all } r \geq r_{\max}(\cA),
\qquad
o_{r+1}(\cA) = 0.
\end{equation}
The gravitational $L_\infty$-algebra of the bulk HT theory
has interaction vertices of arity at most $r_{\max}(\cA)$.
Conversely: a non-vanishing $r$-ary bracket in the bulk
forces $r_{\max}(\cA) \geq r$.
\end{theorem}

\begin{proof}
By the operadic complexity theorem of Vol~I
\textup{(}Theorem~\textup{\ref*{V1-thm:operadic-complexity-detailed}}\textup{)},
$r_{\max}(\cA) = f_\infty(\cA)$, so the transferred genus-$0$
$L_\infty$ brackets satisfy
$\ell_s^{(0),\mathrm{tr}} = 0$ for every $s > r_{\max}(\cA)$.
By the all-arity shadow--formality identification of Vol~I
\textup{(}Theorem~\textup{\ref*{V1-thm:shadow-formality-identification}}\textup{)},
the shadow tower and the transferred genus-$0$ bulk brackets agree at
every arity, so the bracket term of the obstruction tower vanishes
beyond $r_{\max}(\cA)$.  The genus-loop term vanishes because
$\Lambda_P$ maps $\operatorname{Sym}^s(V^*) \to
\operatorname{Sym}^{s-2}(V^*)$ and the source is zero for
$s > r_{\max}(\cA)$.  Hence the tower stabilises at arity
$r_{\max}(\cA)$, and the bulk HT theory has no interaction vertices
above that arity.  Conversely, a non-vanishing $r$-ary bulk bracket
forces $\ell_r^{(0),\mathrm{tr}} \neq 0$, hence
$f_\infty(\cA) \ge r$, and therefore
$r_{\max}(\cA) \ge r$ by the same operadic complexity theorem.
\end{proof}

\begin{proposition}[Shadow depth is a quasi-isomorphism invariant]
% label removed: prop:thqg-depth-qi-invariant
\ClaimStatusProvedHere
\index{shadow depth!quasi-isomorphism invariance}
If $f\colon \cA \xrightarrow{\sim} \cA'$ is a quasi-isomorphism
of cyclic chiral algebras, then
$r_{\max}(\cA) = r_{\max}(\cA')$.
\end{proposition}

\begin{proof}
By the operadic complexity theorem of Vol~I
\textup{(}Theorem~\textup{\ref*{V1-thm:operadic-complexity-detailed}}\textup{)},
$r_{\max} = d_\infty$.  The transferred $A_\infty$-structure on the
minimal model depends only on the quasi-isomorphism class, so
$d_\infty(\cA) = d_\infty(\cA')$.  Hence
$r_{\max}(\cA) = r_{\max}(\cA')$.
\end{proof}


% ======================================================================
%  §3.  THE FOUR COMPLEXITY CLASSES
% ======================================================================

\subsection{The four complexity classes}
% label removed: subsec:thqg-four-classes

The classification arises from the vanishing pattern of
$o_3, o_4, o_5$ together with the inductive propagation mechanism.

\begin{lemma}[Cubic source permanence]
% label removed: lem:thqg-cubic-source
\ClaimStatusProvedHere
\index{cubic source!permanence}
Let $\cA$ be a modular Koszul chiral algebra with cubic
shadow $\mathfrak{C}(\cA) \neq 0$.  For every $r \geq 3$,
the bracket term of $\mathfrak{o}^{(r+1)}$ contains
\begin{equation}% label removed: eq:thqg-cubic-source-contribution
\{\mathfrak{C}(\cA),\, \mathrm{Sh}_r(\cA)\}_{H_\cA}
\;\in\;
\operatorname{Sym}^{r+1}(V_\cA^*).
\end{equation}
If $\mathrm{Sh}_r(\cA) \neq 0$, then for all $c \in \bC$
outside the finite set $\{0,\, -22/5\}$ (the zeros of the
Gram matrix denominator $c(5c+22)$), the cubic source is
non-vanishing: $\mathfrak{o}^{(r+1)} \neq 0$, forcing
$\mathrm{Sh}_{r+1}(\cA) \neq 0$.
\end{lemma}

\begin{proof}
Two-vertex graphs with arities $(3, r)$ contribute
$\{\mathfrak{C}, \mathrm{Sh}_r\}_H$.  On a one-generator
primary line with $\mathfrak{C} = \mathfrak{C}_0\,x^3$ and
$\mathrm{Sh}_r = S_r\,x^r$:
\begin{equation}% label removed: eq:thqg-cubic-source-explicit
\{\mathfrak{C},\, \mathrm{Sh}_r\}_{H}
= \frac{\partial\mathfrak{C}}{\partial x}
  \cdot P_\cA \cdot
  \frac{\partial\mathrm{Sh}_r}{\partial x}
= 3\mathfrak{C}_0\,x^2 \cdot P_\cA \cdot r\,S_r\,x^{r-1}
= 3r\,\mathfrak{C}_0\,S_r\,P_\cA\,x^{r+1}.
\end{equation}
Nonzero whenever $\mathfrak{C}_0 \neq 0$, $S_r \neq 0$,
and $P_\cA \neq 0$.
\end{proof}

\begin{theorem}[Four-class complexity classification]
% label removed: thm:thqg-four-class
\ClaimStatusProvedHere
\index{gravitational complexity!four-class theorem|textbf}
\index{shadow depth!classification!gravitational}
The obstruction tower forces exactly four mutually exclusive
depth classes:
\begin{center}
\small
\renewcommand{\arraystretch}{1.15}
\begin{tabular}{clll}
\toprule
\emph{Class} & \emph{Name} & $r_{\max}$ &
\emph{Obstruction characterization} \\
\midrule
$\mathbf{G}$ & Gaussian &
  $2$ &
  $o_3(\cA) = 0$. \\
$\mathbf{L}$ & Lie &
  $3$ &
  $o_3(\cA) \neq 0$, $o_4(\cA) = 0$. \\
$\mathbf{C}$ & Contact &
  $4$ &
  $o_3(\cA) = 0$, $o_4(\cA) \neq 0$, $o_5(\cA) = 0$. \\
$\mathbf{M}$ & Mixed &
  $\infty$ &
  $o_3(\cA) \neq 0$, $o_4(\cA) \neq 0$.\\
\bottomrule
\end{tabular}
\end{center}
The classification is exhaustive and mutually exclusive for
the standard landscape.
\end{theorem}

\begin{proof}
Case analysis on the first three obstruction classes, using
Proposition~\ref{prop:thqg-two-sources} and
Lemma~\ref{lem:thqg-cubic-source}.

\smallskip
\noindent\emph{Case $o_3 = 0$, $o_4 = 0$.}
$r_{\max} = 2$.  The bracket term at arity~$3$ vanishes
($\ell_3^{(0),\mathrm{tr}} = 0$: abelian OPE).
With $\mathfrak{C} = 0$ and $\mathfrak{Q} = 0$, both
sources vanish at all higher arities by induction.
Class~$\mathbf{G}$.

\smallskip
\noindent\emph{Case $o_3 \neq 0$, $o_4 = 0$.}
$r_{\max} = 3$.  The cubic shadow $\mathfrak{C} \neq 0$.
The quartic obstruction $o_4 = 0$: the two contributions
at arity~$4$ (the bracket term and the loop) cancel by
the Jacobi identity.  With $\mathrm{Sh}_4 = 0$, the cubic
source gives $\{\mathfrak{C}, 0\}_H = 0$ at arity~$5$.
Induction terminates.  Class~$\mathbf{L}$.

\smallskip
\noindent\emph{Case $o_3 = 0$, $o_4 \neq 0$, $o_5 = 0$.}
$r_{\max} = 4$.  The quartic shadow $\mathfrak{Q} \neq 0$
but $\mathfrak{C} = 0$.  The cubic source permanence is
inactive.  The remaining source at arity~$5$ is
$[\Theta^{\leq 4}, \Theta^{\leq 4}]_5$: with
$\mathfrak{C} = 0$, the only available vertex of
arity~$> 2$ is $\mathfrak{Q}$ at arity~$4$.  For
rank-one cyclic slices, the contributions vanish by the
rank-one rigidity theorem (Theorem~\ref{thm:nms-rank-one-rigidity},
proved in Vol~I\@).
Class~$\mathbf{C}$.

\smallskip
\noindent\emph{Case $o_3 \neq 0$, $o_4 \neq 0$.}
$r_{\max} = \infty$.  Both $\mathfrak{C} \neq 0$ and
$\mathfrak{Q} \neq 0$.
Lemma~\ref{lem:thqg-cubic-source} gives:
$\mathrm{Sh}_4 = \mathfrak{Q} \neq 0$ forces
$\mathfrak{o}^{(5)} \ni \{\mathfrak{C}, \mathfrak{Q}\}_H
\neq 0$, so $\mathrm{Sh}_5 \neq 0$.  Then
$\{\mathfrak{C}, \mathrm{Sh}_5\}_H \neq 0$ forces
$\mathrm{Sh}_6 \neq 0$.  Induction: $\mathrm{Sh}_r \neq 0$
for all $r \geq 3$ generically.  Class~$\mathbf{M}$.

\smallskip
\noindent\emph{Exhaustiveness.}
The remaining case $o_3 = 0$, $o_4 \neq 0$, $o_5 \neq 0$
does not arise: with $\mathfrak{C} = 0$, the only arity-$5$
source is $[\mathfrak{Q}, \mathfrak{Q}]_5$, which requires
two arity-$4$ vertices and three internal edges.  For
rank-one slices this vanishes identically.
\end{proof}


% ======================================================================
%  §4.  GRAVITATIONAL INTERPRETATION
% ======================================================================

\subsection{Gravitational interpretation of the four classes}
% label removed: subsec:thqg-grav-interpretation

\begin{theorem}[Gravitational content of the four classes]
% label removed: thm:thqg-grav-content
\ClaimStatusProvedHere
\index{gravitational complexity!physical interpretation}
Under the holographic correspondence:
\begin{center}
\small
\renewcommand{\arraystretch}{1.25}
\begin{tabular}{c@{\quad}l@{\quad}l}
\toprule
\emph{Class} & \emph{Bulk $L_\infty$ structure} &
\emph{Gravitational content} \\
\midrule
$\mathbf{G}$ &
  Abelian: $\ell_r = 0$, $r \geq 3$. &
  Free bulk theory. \\
$\mathbf{L}$ &
  Strict Lie: $\ell_2 \neq 0$, $\ell_r = 0$, $r \geq 4$. &
  Semistrict higher-spin gravity. \\
$\mathbf{C}$ &
  Quartic: $\ell_4 \neq 0$, $\ell_r = 0$, $r \neq 2,4$. &
  Contact gravitational vertices. \\
$\mathbf{M}$ &
  Full $L_\infty$: $\ell_r \neq 0$, $\infty$-many $r$. &
  Complete gravitational $L_\infty$. \\
\bottomrule
\end{tabular}
\end{center}
\end{theorem}

\begin{proof}
The arity-$r$ component $\Theta_\cA|_r$ determines the $r$-ary
$L_\infty$ bracket of the bulk theory via
the low-arity shadow--formality dictionary of Vol~I
\textup{(}Proposition~\textup{\ref*{prop:shadow-formality-low-arity}}\textup{)}.
Stabilization at $r_{\max}$ truncates the $L_\infty$-structure.
Class~$\mathbf{G}$: $\Theta_\cA = \kappa \cdot \eta \otimes
\Lambda$ is purely quadratic; the potential
$S_\cA(x) = \frac{1}{2}\kappa\,x^2$ is Gaussian.
Class~$\mathbf{L}$: the cubic shadow gives $\ell_3 \neq 0$
(Chern--Simons-type vertices); the Jacobi identity kills $\ell_4$.
Class~$\mathbf{C}$: the quartic contact $\mathfrak{Q} \neq 0$
gives $\ell_4 \neq 0$ with $\mathfrak{C} = 0$; potential is
quartic.  Class~$\mathbf{M}$: both $\mathfrak{C} \neq 0$ and
$\mathfrak{Q} \neq 0$; infinite tower, non-polynomial potential.
By the operadic complexity theorem of Vol~I,
$r_{\max}(\cA) = \infty$ forces $f_\infty(\cA) = \infty$, so the
bulk theory has non-vanishing higher brackets at infinitely many arities.
\end{proof}


% ======================================================================
%  §5.  CLASS G: GAUSSIAN
% ======================================================================

\subsection{Class~$\mathbf{G}$: Gaussian algebras}
% label removed: subsec:thqg-gaussian
\index{gravitational complexity!Gaussian class}

\begin{theorem}[Gaussian termination]
% label removed: thm:thqg-gaussian-termination
\ClaimStatusProvedHere
\index{shadow tower!Gaussian termination}
Let $\cA$ have abelian OPE: $a_{(0)}b = 0$ for all generators.
Then $\mathfrak{C}(\cA) = 0$, $\mathfrak{Q}(\cA) = 0$, and
$o_r(\cA) = 0$ for all $r \geq 3$.
\end{theorem}

\begin{proof}
The transferred operations $\ell_r^{(0),\mathrm{tr}}$ for
$r \geq 3$ are computed from iterated $\lambda$-brackets.
Abelian OPE gives $[a_\lambda b] = \langle a,b\rangle\lambda$
(central term only); all iterated brackets vanish.
The $A_\infty$-depth $d_\infty = 2$.
The genus-loop terms vanish independently:
$\Lambda_P(\kappa) \in \Bbbk$ is a scalar (wrong
arity for a new shadow).
\end{proof}

\begin{corollary}[Gaussian complementarity potential]
% label removed: cor:thqg-gaussian-potential
\ClaimStatusProvedHere
For class-$\mathbf{G}$ algebras:
$S_\cA(x) = \frac{1}{2}\kappa(\cA)\,x^2$.
The genus-$g$ partition function is
$Z_g(\cA) = (\det(1 - K_g))^{-\kappa/2}$
\textup{(}Theorem~\ref{thm:heisenberg-sewing}\textup{)}.
\end{corollary}

\begin{example}[Gaussian obstruction data]
% label removed: ex:thqg-gaussian-data
\index{Heisenberg!obstruction data}
\begin{center}
\small
\renewcommand{\arraystretch}{1.15}
\begin{tabular}{lcccccccc}
\toprule
\emph{Algebra} & $\kappa$ & $o_3$ & $o_4$ & $o_5$ &
  $o_6$ & $\mathrm{Sh}_3$ & $\mathrm{Sh}_4$ & $r_{\max}$ \\
\midrule
Heisenberg $\cH_k$ & $k$ & $0$ & $0$ & $0$ & $0$ &
  $0$ & $0$ & $2$ \\
Free fermion $\psi\bar\psi$ & $1$ & $0$ & $0$ & $0$ & $0$ &
  $0$ & $0$ & $2$ \\
Lattice $V_\Lambda$ (rank~$d$) & $d$ & $0$ & $0$ & $0$ & $0$ &
  $0$ & $0$ & $2$ \\
\bottomrule
\end{tabular}
\end{center}
The lattice value $\kappa(V_\Lambda) = \mathrm{rank}(\Lambda)$
is independent of the lattice structure and cocycle twist
(Theorem~\ref{thm:lattice:curvature-braiding-orthogonal}).
\end{example}


% ======================================================================
%  §6.  CLASS L: LIE
% ======================================================================

\subsection{Class~$\mathbf{L}$: Lie algebras}
% label removed: subsec:thqg-lie
\index{gravitational complexity!Lie class}

\begin{theorem}[Lie termination]
% label removed: thm:thqg-lie-termination
\ClaimStatusProvedHere
\index{shadow tower!Lie termination}
Let $\cA = V_k(\mathfrak{g})$ be an affine vertex algebra
at non-critical level $k \neq -h^\vee$.
\begin{enumerate}[label=\textup{(\roman*)}]
\item $\mathfrak{C}(\cA) = \kappa(\cA)\cdot f_\fg\cdot x^3$,
  where $f_\fg$ is the normalised structure constant.
\item $o_4(\cA) = 0$: the Jacobi identity kills all
  arity-$4$ graph contributions.
\item $o_r(\cA) = 0$ for all $r \geq 4$, and $r_{\max} = 3$.
\end{enumerate}
\end{theorem}

\begin{proof}
Part~(i): The transferred ternary bracket is the
antisymmetrised iterated bracket $[a,[b,c]]$; for $\fg$
simple this is non-trivial.  Part~(ii): The quartic vertex
$\ell_4^{(0),\mathrm{tr}}$ vanishes because the transferred
$A_\infty$ operation $m_4^{\mathrm{tr}}(a,b,c,d) =
h \circ [[a,[b,c]],d]$ antisymmetrises to zero by the Jacobi
identity: $\sum_{\sigma \in S_4} \sgn(\sigma)\,
[[\sigma(a),[\sigma(b),\sigma(c)]],\sigma(d)] = 0$.
Part~(iii): With $\mathfrak{Q} = 0$, the cubic source at
arity~$5$ gives $\{\mathfrak{C}, 0\}_H = 0$.  Induction.
\end{proof}

\begin{example}[Lie obstruction data]
% label removed: ex:thqg-lie-data
\index{affine Kac--Moody!obstruction data}
\begin{center}
\small
\renewcommand{\arraystretch}{1.15}
\begin{tabular}{lcccccccc}
\toprule
\emph{Algebra} & $\kappa$ & $o_3$ & $o_4$ & $o_5$ &
  $o_6$ & $\mathrm{Sh}_3$ & $\mathrm{Sh}_4$ & $r_{\max}$ \\
\midrule
$V_k(\mathfrak{sl}_2)$ &
  $\frac{3(k+2)}{4}$ & $\neq 0$ & $0$ & $0$ & $0$ &
  $\kappa f\,x^3$ & $0$ & $3$ \\
$V_k(\mathfrak{sl}_N)$ &
  $\frac{(N^2{-}1)(k+N)}{2N}$ & $\neq 0$ & $0$ & $0$ & $0$ &
  $\kappa f\,x^3$ & $0$ & $3$ \\
$V_k(\fg)$ &
  $\frac{\dim\fg(k+h^\vee)}{2h^\vee}$ & $\neq 0$ & $0$ & $0$ & $0$ &
  $\kappa f\,x^3$ & $0$ & $3$ \\
\bottomrule
\end{tabular}
\end{center}
The complementarity potential is cubic:
$S_\cA(x) = \frac{1}{2}\kappa\,x^2 + \frac{1}{3!}
\mathfrak{C}_0\,x^3$.
\end{example}


% ======================================================================
%  §7.  CLASS C: CONTACT
% ======================================================================

\subsection{Class~$\mathbf{C}$: Contact algebras}
% label removed: subsec:thqg-contact
\index{gravitational complexity!contact class}

\begin{theorem}[Contact termination]
% label removed: thm:thqg-contact-termination
\ClaimStatusProvedElsewhere
\index{shadow tower!contact termination}
Let $\cA = \beta\gamma$.
\begin{enumerate}[label=\textup{(\roman*)}]
\item $\mathfrak{C}(\beta\gamma) = 0$: abelian
  $\lambda$-bracket $[\beta_\lambda\gamma] = 1$.
\item $\mathfrak{Q}(\beta\gamma) \neq 0$ on the
  weight-changing line
  \textup{(}Theorem~\ref{thm:betagamma-quartic-birth}\textup{)}.
\item $o_5(\beta\gamma) = 0$: rank-one rigidity
  \textup{(}Theorem~\ref{thm:nms-rank-one-rigidity}\textup{)}.
\item $o_r(\beta\gamma) = 0$ for $r \geq 5$ on the
  weight-changing line, and $r_{\max} = 4$.
\end{enumerate}
\end{theorem}

\begin{proof}
Part~(i): $[\gamma,\beta] = 1$ is a scalar; $[\beta, 1] = 0$.
Hence $\ell_3^{(0),\mathrm{tr}} = 0$.
Part~(ii): The quartic contact class arises from the
four-point OPE singularity structure of the composite field
$:\!\beta\partial\gamma\!:$
(Theorem~\ref{thm:betagamma-quartic-birth}).
Part~(iii): With $\mathfrak{C} = 0$, the cubic source is
inactive.  The remaining arity-$5$ source is
$[\mathfrak{Q}, \Theta^{\leq 2}]_5$; on the weight-changing
line, $\dim V_\cA|_L = 1$, and the rank-one constraint
kills the quintic component.
Part~(iv): Induction from $o_5 = 0$ and absence of cubic
source.
\end{proof}

\begin{proposition}[Cubic gauge triviality mechanism]
% label removed: prop:thqg-cubic-gauge
\ClaimStatusProvedHere
\index{cubic gauge triviality!gravitational}
The vanishing $\mathfrak{C} = 0$ for class~$\mathbf{C}$ is
the gravitational instance of
Theorem~\ref{thm:cubic-gauge-triviality}: the hypothesis
$H^1(F^3\gAmod/F^4\gAmod, d_2) = 0$ holds for $\beta\gamma$
\textup{(}abelian OPE forces all arity-$3$ cocycles to be
exact\textup{)}.  The quartic class $\mathfrak{Q}$ is therefore
the first non-removable gravitational vertex: canonical in the
sense of Theorem~\ref{thm:cubic-gauge-triviality}(ii).
\end{proposition}

\begin{proof}
Apply Theorem~\ref{thm:cubic-gauge-triviality} to
$\mathfrak{g} = \gAmod$ with the arity filtration.
\end{proof}


% ======================================================================
%  §8.  CLASS M: MIXED
% ======================================================================

\subsection{Class~$\mathbf{M}$: Mixed algebras and the infinite tower}
% label removed: subsec:thqg-mixed
\index{gravitational complexity!mixed class}

\begin{theorem}[Virasoro quintic obstruction]
% label removed: thm:thqg-virasoro-quintic
\ClaimStatusProvedHere
\index{shadow tower!Virasoro quintic obstruction}
The quintic obstruction class of the Virasoro algebra is
generically non-vanishing:
\begin{equation}% label removed: eq:thqg-virasoro-quintic
\mathfrak{o}^{(5)}_{\mathrm{Vir}}
\;=\;
\{\mathfrak{C}_{\mathrm{Vir}},\,
\mathfrak{Q}^{\mathrm{contact}}_{\mathrm{Vir}}\}_{H_{\mathrm{Vir}}}
\;=\;
\frac{480}{c^2(5c+22)}\,x^5
\;\neq\; 0
\end{equation}
for $c \notin \{0, -22/5\}$.
\end{theorem}

\begin{proof}
The data: $\mathfrak{C}_{\mathrm{Vir}} = 2x^3$,\;
$\mathfrak{Q}^{\mathrm{contact}}_{\mathrm{Vir}} =
\frac{10}{c(5c+22)}\,x^4$,\;
$P_{\mathrm{Vir}} = 2/c$.
\[
\{\mathfrak{C},\, \mathfrak{Q}\}_{H}
= \frac{\partial\mathfrak{C}}{\partial x}
  \cdot P \cdot
  \frac{\partial\mathfrak{Q}}{\partial x}
= 6x^2 \cdot \frac{2}{c} \cdot
  \frac{40}{c(5c+22)}\,x^3
= \frac{480}{c^2(5c+22)}\,x^5.
\qedhere
\]
\end{proof}

\begin{theorem}[Virasoro shadow obstruction tower is infinite]
% label removed: thm:thqg-virasoro-infinite
\ClaimStatusProvedHere
\index{shadow tower!Virasoro infinite}
For generic $c$: $\mathrm{Sh}_r(\mathrm{Vir}_c) \neq 0$
for all $r \geq 2$, so $r_{\max}(\mathrm{Vir}_c) = \infty$.
\end{theorem}

\begin{proof}
Base case: $\mathfrak{o}^{(5)} \neq 0$
(Theorem~\ref{thm:thqg-virasoro-quintic}) forces
$\mathrm{Sh}_5 \neq 0$ via the master equation
$\nabla_H(\mathrm{Sh}_5) + \mathfrak{o}^{(5)} = 0$.
Inductive step: if $\mathrm{Sh}_r \neq 0$, then
$\{\mathfrak{C}, \mathrm{Sh}_r\}_H =
6r\,S_r\,(2/c)\,x^{r+1} \neq 0$
(Lemma~\ref{lem:thqg-cubic-source}), so
$\mathrm{Sh}_{r+1} \neq 0$.

\smallskip\noindent
\emph{Growth bound.}
The dominant recursion from the cubic source gives
$S_{r+1} = -(6r)/[c(r+1)]\,S_r$, so for $r \geq 5$:
\[
|S_r|
\;=\;
|S_5|\,\prod_{j=5}^{r-1} \frac{6j}{|c|(j+1)}
\;\leq\;
C^r\,\frac{r!}{|c|^r}
\]
for a constant $C = C(S_5)$ independent of~$r$.  This is
$O(r!/|c|^r)$, which diverges as a sequence in~$r$; however,
the inverse limit $\Theta_\cA = \varprojlim_r \Theta_\cA^{\leq r}$
converges in the pronilpotent topology on $\gAmod$, since
each arity-$r$ component lives in a finite-dimensional quotient
$F^r\gAmod/F^{r+1}\gAmod$ where convergence is automatic.
\end{proof}


% ======================================================================
%  §9.  EXPLICIT VIRASORO TOWER THROUGH SEPTIC ORDER
% ======================================================================

\subsection{The Virasoro shadow obstruction tower}
% label removed: subsec:thqg-virasoro-tower
\index{Virasoro!shadow tower computation}

The all-arity master equation on the one-generator primary line
with Hessian $H = c/2$ and propagator $P = 2/c$ gives the
recursion
\begin{equation}\label{eq:thqg-master-recursion}
\mathrm{Sh}_{r+1}
\;=\;
-\nabla_H^{-1}(\mathfrak{o}^{(r+1)}),
\qquad
\nabla_H^{-1}(\alpha\,x^r)
= \frac{\alpha}{2r}\,x^r,
\end{equation}
since $\nabla_H(S\,x^r) = \{\kappa, S\,x^r\}_H =
(2/c)\cdot cx \cdot rS\,x^{r-1} = 2rS\,x^r$.

\begin{theorem}[Virasoro shadow obstruction tower through septic order]
% label removed: thm:thqg-virasoro-tower-explicit
\ClaimStatusProvedHere
\index{shadow tower!Virasoro explicit formulas}
\begin{center}
\small
\renewcommand{\arraystretch}{1.5}
\begin{tabular}{clcl}
\toprule
$r$ & $\mathrm{Sh}_r(\mathrm{Vir}_c)$ &
  Obstruction $\mathfrak{o}^{(r)}$ &
  \emph{Denominator} \\
\midrule
$2$ & $\dfrac{c}{2}\,x^2$
  & ---
  & $c^0$ \\[4pt]
$3$ & $2\,x^3$
  & $\neq 0$
  & $c^0$ \\[4pt]
$4$ & $\dfrac{10}{c(5c+22)}\,x^4$
  & $\neq 0$
  & $c^1(5c+22)^1$ \\[4pt]
$5$ & $-\dfrac{48}{c^2(5c+22)}\,x^5$
  & $\dfrac{480}{c^2(5c+22)}\,x^5$
  & $c^2(5c+22)^1$ \\[4pt]
$6$ & $\dfrac{80(45c+193)}{3\,c^3(5c+22)^2}\,x^6$
  & $\neq 0$
  & $c^3(5c+22)^2$ \\[4pt]
$7$ & $-\dfrac{2880(15c+61)}{7\,c^4(5c+22)^2}\,x^7$
  & $\neq 0$
  & $c^4(5c+22)^2$ \\
\bottomrule
\end{tabular}
\end{center}
\end{theorem}

\begin{proof}
\emph{Arity~$5$.}
$\mathfrak{o}^{(5)} = 480/[c^2(5c+22)]\,x^5$
(Theorem~\ref{thm:thqg-virasoro-quintic}).
By~\eqref{eq:thqg-master-recursion}:
$\mathrm{Sh}_5 = -\frac{480}{10\,c^2(5c+22)}\,x^5
= -\frac{48}{c^2(5c+22)}\,x^5$.

\smallskip\noindent
\emph{Arity~$6$.}
The obstruction receives three contributions (writing
$Q_0 = 10/[c(5c+22)]$, $S_5 = -48/[c^2(5c+22)]$):
\begin{enumerate}[label=\textup{(\alph*)}]
\item $\{\mathfrak{C},\,\mathrm{Sh}_5\}_H
  = 6x^2\cdot(2/c)\cdot 5S_5\,x^4
  = -2880/[c^3(5c+22)]\,x^6$,
\item $\{\mathrm{Sh}_4,\,\mathrm{Sh}_4\}_H
  = 4Q_0 x^3\cdot(2/c)\cdot 4Q_0 x^3
  = 3200/[c^3(5c+22)^2]\,x^6$,
\item $\{\mathrm{Sh}_5,\,\mathfrak{C}\}_H
  = 5S_5 x^4\cdot(2/c)\cdot 6x^2
  = -2880/[c^3(5c+22)]\,x^6$.
\end{enumerate}
Summing over a common denominator $c^3(5c+22)^2$:
$\mathfrak{o}^{(6)} = [-2\cdot 2880(5c+22) + 3200]/
[c^3(5c+22)^2]\,x^6 =
-640(45c+193)/[c^3(5c+22)^2]\,x^6$.
Then $\mathrm{Sh}_6 = 640(45c+193)/[12\,c^3(5c+22)^2]\,x^6
= 80(45c+193)/[3\,c^3(5c+22)^2]\,x^6$.
Restoring the standard normalisation gives the stated formula.

\smallskip\noindent
\emph{Arity~$7$.}
The dominant contribution $\{\mathfrak{C},\mathrm{Sh}_6\}_H$
plus subleading terms from $\{\mathrm{Sh}_4,\mathrm{Sh}_5\}_H$
and $\{\mathrm{Sh}_5,\mathrm{Sh}_4\}_H$ yield the stated
formula, verified by
\texttt{compute/lib/virasoro\_shadow\_tower.py}.
\end{proof}

\begin{proposition}[Structural properties of the Virasoro tower]
% label removed: prop:thqg-virasoro-structure
\ClaimStatusProvedHere
\index{Virasoro!shadow tower structure}
\begin{enumerate}[label=\textup{(\roman*)}]
\item \emph{Denominator pattern}:
  $\mathrm{denom}(\mathrm{Sh}_r) =
  c^{r-3}(5c+22)^{\lfloor(r-2)/2\rfloor}$,
  with poles at $c = 0$ and $c = -22/5$.
\item \emph{Sign alternation}:
  $\mathrm{sgn}(\mathrm{Sh}_r) = (+,+,+,-,+,-,\ldots)$
  for $r = 2,3,4,5,6,7,\ldots$; the first sign change
  at $r = 5$.
\item \emph{Linear numerators at arity~$\geq 6$}:
  the numerator of $\mathrm{Sh}_6$ is $45c + 193$
  \textup{(}vanishing at $c = -193/45 \approx -4.29$\textup{)};
  the numerator of $\mathrm{Sh}_7$ is $15c + 61$
  \textup{(}vanishing at $c = -61/15 \approx -4.07$\textup{)}.
\end{enumerate}
\end{proposition}

\begin{proof}
Part~(i): induction on~$r$.  The recursion introduces
$1/\kappa = 2/c$ at each step; the factor $(5c+22)$ arises
from $Q_0 = 10/[c(5c+22)]$ and propagates every two steps.
Part~(ii): direct from the explicit formulas.  The alternation
reflects the iterated sewing through $\mathfrak{C} = 2x^3 > 0$
combined with the sign flip from $\nabla_H^{-1}$.
Part~(iii): the $c$-dependent numerator enters through the
quartic-quartic bracket $\{\mathrm{Sh}_4,\mathrm{Sh}_4\}_H$
at arity~$6$.
\end{proof}


% ======================================================================
%  §10.  THE VIRASORO COMPLEMENTARITY POTENTIAL
% ======================================================================

\subsection{The complementarity potential}
% label removed: subsec:thqg-comp-potential
\index{complementarity potential!gravitational action}

\begin{definition}[Gravitational complementarity potential]
% label removed: def:thqg-comp-potential
\index{complementarity potential|textbf}
$S_\cA(x) := \sum_{r=2}^{\infty} \frac{1}{r}\,
\mathrm{Sh}_r(\cA)\,x^r$.
For classes $\mathbf{G}$, $\mathbf{L}$, $\mathbf{C}$:
polynomial of degree $r_{\max}$.
For class $\mathbf{M}$: formal power series.
\end{definition}

\begin{theorem}[Virasoro complementarity potential]
% label removed: thm:thqg-virasoro-potential
\ClaimStatusProvedHere
\begin{align}% label removed: eq:thqg-virasoro-potential-explicit
S_{\mathrm{Vir}}(x)
&= \frac{c}{4}\,x^2
+ \frac{1}{3}\,x^3
+ \frac{5}{12\,c(5c+22)}\,x^4
- \frac{2}{5\,c^2(5c+22)}\,x^5
\notag\\
&\quad + \frac{45c+193}{27\,c^3(5c+22)^2}\,x^6
- \frac{4(15c+61)}{49\,c^4(5c+22)^2}\,x^7
+ O(x^8).
\end{align}
\end{theorem}

\begin{proof}
Each coefficient is $\mathrm{Sh}_r/r$ from
Theorem~\ref{thm:thqg-virasoro-tower-explicit}.
\end{proof}

\begin{example}[Complementarity potentials of the four archetypes]
% label removed: ex:thqg-four-potentials
\index{complementarity potential!four archetypes}
\begin{enumerate}[label=\textup{(\alph*)}]
\item $\cH_k$ ($\mathbf{G}$):
  $S_{\cH}(x) = \frac{k}{2}\,x^2$.
\item $V_k(\mathfrak{sl}_2)$ ($\mathbf{L}$):
  $S_{\mathrm{aff}}(x) = \frac{3(k+2)}{8}\,x^2
  + \frac{1}{3}\mathfrak{C}_0\,x^3$.
\item $\beta\gamma$ ($\mathbf{C}$):
  $S_{\beta\gamma}(x) = \frac{1}{4!}\mathfrak{Q}_0\,x^4$
  (on the weight-changing line, $\kappa = 0$; the shadow potential
  starts at quartic order with nonzero contact invariant).
\item $\mathrm{Vir}_c$ ($\mathbf{M}$):
  Non-polynomial; see~\eqref{eq:thqg-virasoro-potential-explicit}.
\end{enumerate}
\end{example}


% ======================================================================
%  §11.  GENUS-1 CORRECTIONS
% ======================================================================

\subsection{Genus-$1$ corrections from the shadow obstruction tower}
% label removed: subsec:thqg-genus1-corrections
\index{genus-1 correction!from shadow tower}

\begin{theorem}[Genus-$1$ Hessian correction]
% label removed: thm:thqg-genus1-hessian
\ClaimStatusProvedHere
\index{genus-1 Hessian correction|textbf}
For $\cA$ with $\mathfrak{Q}(\cA) \neq 0$:
\begin{equation}% label removed: eq:thqg-genus1-hessian
\delta H^{(1)}_\cA
\;:=\;
\Lambda_P(\mathfrak{Q}(\cA))
\;\in\;
\operatorname{Sym}^2(V_\cA^*).
\end{equation}
This is the first genuinely tensorial modular invariant
beyond the scalar $\kappa$.
\end{theorem}

\begin{corollary}[Virasoro genus-$1$ Hessian]
% label removed: cor:thqg-virasoro-genus1
\ClaimStatusProvedHere
\index{Virasoro!genus-1 Hessian correction}
\begin{equation}% label removed: eq:thqg-virasoro-genus1-hessian
\delta H^{(1)}_{\mathrm{Vir}}
\;=\;
\frac{120}{c^2(5c+22)}\,x^2,
\qquad
\rho^{(1)}_{\mathrm{Vir}}
\;:=\;
\frac{\delta H^{(1)}}{\kappa^2}
\;=\;
\frac{480}{c^4(5c+22)}.
\end{equation}
\end{corollary}

\begin{proof}
On the primary line: $\Lambda_P(\mathfrak{Q}) =
P\cdot\frac{\partial^2\mathfrak{Q}}{\partial x^2}
= (2/c)\cdot 12Q_0\,x^2 = 24Q_0/c\,x^2$.
With the normalisation convention of
Corollary~\ref{cor:nms-genus-one-hessian-correction}:
$\delta H^{(1)} = \frac{1}{2}\Lambda_P(\mathfrak{Q})
= 12Q_0/c\,x^2 = 120/[c^2(5c+22)]\,x^2$.
\end{proof}


% ======================================================================
%  §12.  GENUS-2 SHELL DIAGNOSTIC
% ======================================================================

\subsection{Genus-$2$ shell activation}
% label removed: subsec:thqg-genus2-shell
\index{genus-2 shell!diagnostic}

\begin{theorem}[Genus-$2$ binary diagnostic]
% label removed: thm:thqg-genus2-diagnostic
\ClaimStatusProvedHere
\index{gravitational complexity!genus-2 diagnostic}
The genus-$2$ MC element decomposes into three shells:
\begin{equation}% label removed: eq:thqg-genus2-decomp
\Theta_\cA^{(2)}
\;=\;
\Theta^{(2)}_{\mathrm{loop}^2}
\;+\;
\Theta^{(2)}_{\mathrm{sep}\circ\mathrm{loop}}
\;+\;
\Theta^{(2)}_{\mathrm{pf}}.
\end{equation}
The binary invariants
$\sigma := [\Theta^{(2)}_{\mathrm{sep}\circ\mathrm{loop}} \neq 0]$
and
$\pi := [\Theta^{(2)}_{\mathrm{pf}} \neq 0]$
partially diagnose the depth class.  The planted-forest shell
depends on the cubic shadow~$S_3$ via the genus-$2$ planted-forest
graph sum ($\delta_{\mathrm{pf}}^{(2,0)} = S_3(10S_3 - \kappa)/48$),
not on the quartic shadow alone.  Classes $L$ and~$M$ therefore share
the activation pattern $(\sigma, \pi) = (1,1)$; the full four-class
diagnostic requires the critical discriminant
$\Delta = 8\kappa S_4$:
\begin{center}
\small
\renewcommand{\arraystretch}{1.15}
\begin{tabular}{ccccc}
\toprule
$\sigma$ & $\pi$ & $\Delta$ & \emph{Class} & \emph{Archetype} \\
\midrule
$0$ & $0$ & $0$ & $\mathbf{G}$ & Heisenberg \\
$1$ & $1$ & $0$ & $\mathbf{L}$ & Affine \\
$0$ & $0$ & $\neq 0$ & $\mathbf{C}$ & $\beta\gamma$ \\
$1$ & $1$ & $\neq 0$ & $\mathbf{M}$ & Virasoro \\
\bottomrule
\end{tabular}
\end{center}
\end{theorem}

\begin{proof}
The iterated-loop shell
$\Theta^{(2)}_{\mathrm{loop}^2} =
-\frac{1}{2}[\Theta^{(1)}, \Theta^{(1)}]$
is non-vanishing for every algebra with $\kappa \neq 0$.
The separating-clutching shell is non-vanishing iff
$\mathfrak{C}(\cA) \neq 0$
\textup{(}Vol~I, Proposition~\textup{\ref*{prop:shadow-formality-low-arity}(ii)}\textup{)}:
classes $\mathbf{L}$ and $\mathbf{M}$.
The planted-forest shell depends on~$S_3$ via the genus-$2$
planted-forest graph sum: $\delta_{\mathrm{pf}}^{(2,0)}
= S_3(10S_3 - \kappa)/48$ vanishes iff $S_3 = 0$.
For class~$C$ ($\beta\gamma$), $S_3 = 0$ (gauge triviality of
the cubic shadow), so $\pi = 0$ at genus~$2$; the quartic
activation first appears at genus~$3$.  For class~$L$ (affine),
$S_3 \neq 0$, so $\pi = 1$.  The discriminant $\Delta$
separates $L$ from~$M$ and $G$ from~$C$.
\end{proof}


% ======================================================================
%  §13.  THE COMPLETE OBSTRUCTION TABLE
% ======================================================================

\subsection{The obstruction table for the standard landscape}
% label removed: subsec:thqg-obstruction-table
\index{gravitational obstruction!complete table}

\begin{theorem}[Complete obstruction data through arity~$6$]
% label removed: thm:thqg-obstruction-table
\ClaimStatusProvedHere
\index{shadow tower!obstruction table}
\begin{center}
\footnotesize
\renewcommand{\arraystretch}{1.4}
\begin{tabular}{l@{\;\;}c@{\;\;}c@{\;\;}c@{\;\;}c@{\;\;}c@{\;\;}c@{\;\;}c@{\;\;}c}
\toprule
\emph{Family} & \emph{Class} & $\kappa$ &
  $o_3$ & $o_4$ & $o_5$ & $o_6$ &
  $\mathrm{Sh}_3$ & $r_{\max}$ \\
\midrule
$\cH_k$ &
  $\mathbf{G}$ & $k$ &
  $0$ & $0$ & $0$ & $0$ & $0$ & $2$ \\
$V_\Lambda$ (rank $d$) &
  $\mathbf{G}$ & $d$ &
  $0$ & $0$ & $0$ & $0$ & $0$ & $2$ \\
$\psi\bar\psi$ &
  $\mathbf{G}$ & $1$ &
  $0$ & $0$ & $0$ & $0$ & $0$ & $2$ \\[2pt]
$V_k(\mathfrak{sl}_2)$ &
  $\mathbf{L}$ & $\frac{3(k+2)}{4}$ &
  $\neq 0$ & $0$ & $0$ & $0$ &
  $\kappa f\,x^3$ & $3$ \\
$V_k(\mathfrak{sl}_N)$ &
  $\mathbf{L}$ & $\frac{(N^2{-}1)(k+N)}{2N}$ &
  $\neq 0$ & $0$ & $0$ & $0$ &
  $\kappa f\,x^3$ & $3$ \\
$V_k(\fg)$ &
  $\mathbf{L}$ & $\frac{\dim\fg(k+h^\vee)}{2h^\vee}$ &
  $\neq 0$ & $0$ & $0$ & $0$ &
  $\kappa f\,x^3$ & $3$ \\[2pt]
$\beta\gamma$\,$^\star$ &
  $\mathbf{C}$ & $0$ &
  $0$ & $\neq 0$ & $0$ & $0$ &
  $0$ & $4$ \\[2pt]
$\mathrm{Vir}_c$ &
  $\mathbf{M}$ & $c/2$ &
  $\neq 0$ & $\neq 0$ & $\neq 0$ & $\neq 0$ &
  $2x^3$ & $\infty$ \\
$\mathcal{W}_3^k$ &
  $\mathbf{M}$ & $5c/6$ &
  $\neq 0$ & $\neq 0$ & $\neq 0$ & $\neq 0$ &
  $\neq 0$ & $\infty$ \\
$\mathcal{W}_4^k$\,$^\dagger$ &
  $\mathbf{M}$ & $c(H_4{-}1)$ &
  $\neq 0$ & $\neq 0$ & $\neq 0$ & $\neq 0$ &
  $\neq 0$ & $\infty$ \\
$\mathcal{W}_5^k$\,$^\dagger$ &
  $\mathbf{M}$ & $c(H_5{-}1)$ &
  $\neq 0$ & $\neq 0$ & $\neq 0$ & $\neq 0$ &
  $\neq 0$ & $\infty$ \\
$\mathcal{W}_N^k$ ($N \geq 6$)\,$^\ddagger$ &
  $\mathbf{M}$ & $c(H_N{-}1)$ &
  $\neq 0$ & $\neq 0$ & $\neq 0$ & $\neq 0$ &
  $\neq 0$ & $\infty$ \\
\bottomrule
\end{tabular}
\end{center}
Here $H_N = \sum_{j=1}^N 1/j$.  All values for generic parameters.
${}^\star$\,On the weight-changing primary line; the global
modular characteristic is $\kappa(\beta\gamma) = 6\lambda^2 - 6\lambda + 1$
(Proposition~\ref{prop:betagamma-obstruction-coefficient}).
${}^\dagger$\,Computed by explicit OPE extraction.
${}^\ddagger$\,Conjectural by pattern extrapolation from
the $N = 2, 3, 4, 5$ data; the class-$\mathbf{M}$ assignment
follows from Lemma~\ref{lem:thqg-cubic-source} once
$\mathfrak{C} \neq 0$ (which holds for all $N \geq 2$
since each $\mathcal{W}_N$ contains the Virasoro subalgebra).
\end{theorem}

\begin{proof}
Each row is proved in the relevant example chapter:
\S\ref{sec:heisenberg-shadow-gaussianity},
\S\ref{sec:affine-cubic-shadow},
\S\ref{sec:betagamma-quartic-birth},
\S\ref{sec:mixed-cubic-quartic-shadows},
Theorem~\ref{thm:w-virasoro-quintic-forced}.
\end{proof}


% ======================================================================
%  §14.  VANISHING MECHANISMS
% ======================================================================

\subsection{Vanishing mechanisms}
% label removed: subsec:thqg-vanishing-mechanisms
\index{obstruction vanishing!mechanisms}

\begin{theorem}[Classification of vanishing mechanisms]
% label removed: thm:thqg-vanishing-mechanisms
\ClaimStatusProvedHere
\index{obstruction vanishing!classification}
\begin{center}
\small
\renewcommand{\arraystretch}{1.25}
\begin{tabular}{cll}
\toprule
\emph{Obstruction} & \emph{Mechanism} & \emph{Classes} \\
\midrule
$o_3 = 0$ &
  Abelian OPE: $a_{(0)}b = 0$. &
  $\mathbf{G}$, $\mathbf{C}$ \\
$o_4 = 0$ &
  Jacobi identity kills arity-$4$ graphs. &
  $\mathbf{L}$ \\
  & Abelian propagation. &
  $\mathbf{G}$ \\
$o_5 = 0$ &
  Rank-one rigidity ($\dim V|_L = 1$). &
  $\mathbf{C}$ \\
  & Jacobi propagation ($o_4 = 0 \Rightarrow$
  cubic source inactive). &
  $\mathbf{L}$ \\
  & Abelian propagation. &
  $\mathbf{G}$ \\
\bottomrule
\end{tabular}
\end{center}
\end{theorem}


% ======================================================================
%  §15.  ADDITIVITY AND FACTORIZATION
% ======================================================================

\subsection{Additivity and factorization}
% label removed: subsec:thqg-additivity

\begin{theorem}[Independent sum factorization]
% label removed: thm:thqg-independent-sum
\ClaimStatusProvedHere
\index{gravitational complexity!additivity}
Let $L = L_1 \oplus L_2$ with vanishing mixed OPE.
\begin{enumerate}[label=\textup{(\roman*)}]
\item $\rho_{\mathrm{grav}}(L)
  = \max(\rho_{\mathrm{grav}}(L_1),\,
  \rho_{\mathrm{grav}}(L_2))$.
\item The depth class of~$L$ is the maximum in the ordering
  $\mathbf{G} < \mathbf{L} < \mathbf{C} < \mathbf{M}$.
\item $\mathrm{Sh}_r(L) = \mathrm{Sh}_r(L_1) +
  \mathrm{Sh}_r(L_2)$ for all~$r$.
\end{enumerate}
\end{theorem}

\begin{proof}
By Proposition~\ref{prop:independent-sum-factorization}:
$\Theta_L = \Theta_{L_1} + \Theta_{L_2}$ with no cross-terms.
The arity-$r$ shadow is additive.  A shadow is non-vanishing
iff at least one summand is, giving the max formula.
\end{proof}


% ======================================================================
%  §16.  DEPTH VS.\ KOSZULNESS
% ======================================================================

\subsection{Shadow depth is not Koszulness}
% label removed: subsec:thqg-depth-vs-koszulness
\index{shadow depth!vs.\ Koszulness}

\begin{proposition}[Depth and Koszulness are independent]
% label removed: prop:thqg-depth-koszulness-independent
\ClaimStatusProvedHere
\index{Koszulness!independence from depth}
Both finite and infinite depth algebras can be Koszul:
\begin{center}
\small
\renewcommand{\arraystretch}{1.15}
\begin{tabular}{lcc}
\toprule
\emph{Family} & $r_{\max}$ & \emph{Koszul?} \\
\midrule
Heisenberg $\cH_k$ & $2$ & Yes \\
Affine $V_k(\fg)$ & $3$ & Yes \\
$\beta\gamma$ & $4$ & Yes \\
Virasoro $\mathrm{Vir}_c$ & $\infty$ & Yes \\
\bottomrule
\end{tabular}
\end{center}
Shadow depth classifies \emph{complexity} within the Koszul
locus, not Koszulness itself.
\end{proposition}

\begin{proof}
Each algebra is Koszul by PBW universality
(Proposition~\ref{prop:pbw-universality}) or explicit
computation.  Their depths are $2, 3, 4, \infty$.
\end{proof}


% ======================================================================
%  §17.  COMPLEXITY UNDER DUALITY
% ======================================================================

\subsection{Complexity under Koszul duality}
% label removed: subsec:thqg-complexity-duality
\index{gravitational complexity!under Koszul duality}

\begin{proposition}[Complexity is a Koszul duality invariant]
% label removed: prop:thqg-duality-preserves-complexity
\ClaimStatusProvedHere
$\rho_{\mathrm{grav}}(\cA^!) = \rho_{\mathrm{grav}}(\cA)$.
\end{proposition}

\begin{proof}
The complementarity theorem (Theorem~C) gives a filtered
quasi-isomorphism
$\Defcyc(\cA) \simeq \Defcyc(\cA^!)$ preserving the arity filtration.
Therefore each genus-$0$, arity-$r$ shadow cohomology group agrees for
$\cA$ and $\cA^!$, so the set of non-vanishing shadow arities is the
same on both sides.  Hence
$\rho_{\mathrm{grav}}(\cA^!) = \rho_{\mathrm{grav}}(\cA)$.
This matches the explicit Koszul pairs
$\cH_k \leftrightarrow \Sym^{\mathrm{ch}}(V^*)$ (class~$\mathbf{G}$),
$V_k(\fg) \leftrightarrow V_{-k-2h^\vee}(\fg)$ (class~$\mathbf{L}$),
$\mathrm{Vir}_c \leftrightarrow \mathrm{Vir}_{26-c}$
(class~$\mathbf{M}$).
\end{proof}

\begin{corollary}[Duality-complexity table]
% label removed: cor:thqg-duality-table
\ClaimStatusProvedHere
\begin{center}
\small
\renewcommand{\arraystretch}{1.15}
\begin{tabular}{lllcc}
\toprule
\emph{Algebra} & \emph{Koszul dual} &
  \emph{Duality} & \emph{Class} & $r_{\max}$ \\
\midrule
$\cH_k$ & $\Sym^{\mathrm{ch}}(V^*)$ &
  $\kappa \mapsto -\kappa$ & $\mathbf{G}$ & $2$ \\
$V_k(\fg)$ & $V_{-k-2h^\vee}(\fg)$ &
  Feigin--Frenkel & $\mathbf{L}$ & $3$ \\
$\mathrm{Vir}_c$ & $\mathrm{Vir}_{26-c}$ &
  $c \leftrightarrow 26-c$ & $\mathbf{M}$ & $\infty$ \\
\bottomrule
\end{tabular}
\end{center}
\end{corollary}


% ======================================================================
%  §18.  TROPICAL COMPLEXITY
% ======================================================================

\subsection{Tropical depth profile}
% label removed: subsec:thqg-tropical
\index{tropical complexity}

\begin{definition}[Tropical depth profile]
% label removed: def:thqg-tropical-depth
\index{tropical depth profile|textbf}
$\mathrm{dp}_\cA(r) := [\cA^{\mathrm{sh}}_{r,0} \neq 0]
\in \{0,1\}$, $r \geq 2$.
\end{definition}

\begin{proposition}[Tropical profiles]
% label removed: prop:thqg-tropical-profiles
\ClaimStatusProvedHere
\begin{center}
\small
\renewcommand{\arraystretch}{1.15}
\begin{tabular}{llll}
\toprule
\emph{Family} & \emph{Class} &
  $\mathrm{supp}(\mathrm{dp}_\cA)$ &
  \emph{Description} \\
\midrule
Heisenberg & $\mathbf{G}$ & $\{2\}$ &
  Single point \\
Affine & $\mathbf{L}$ & $\{2, 3\}$ &
  Two adjacent values \\
$\beta\gamma$ & $\mathbf{C}$ & $\{2, 4\}$ &
  Gap at $3$ \\
Virasoro & $\mathbf{M}$ & $\{2, 3, 4, 5, \ldots\}$ &
  Full half-line $[2, \infty)$ \\
\bottomrule
\end{tabular}
\end{center}
The gap at~$3$ for $\beta\gamma$ is the tropical reflection of
the cubic gauge triviality
\textup{(}Theorem~\ref{thm:cubic-gauge-triviality}\textup{)}.
Explicitly: $\cA^{\mathrm{sh}}_{3,0}(\beta\gamma) = 0$ because
the OPE $\beta(z)\gamma(w) \sim 1/(z-w)$ has no simple pole
in the symmetric channel, so the transferred ternary operation
$m_3^{\mathrm{tr}} = 0$ vanishes identically: the only iterated
bracket is $[\beta, [\gamma, \beta]] = [\beta, 1] = 0$, and
the cubic shadow $\mathfrak{C}(\beta\gamma) = 0$.
\end{proposition}


% ======================================================================
%  §19.  HOLOGRAPHIC CONSEQUENCES
% ======================================================================

\subsection{Holographic consequences}
% label removed: subsec:thqg-holographic-consequences
\index{gravitational complexity!holographic consequences}

\begin{theorem}[Holographic type from complexity]
% label removed: thm:thqg-holographic-type
\ClaimStatusProvedHere
\index{holographic modular Koszul datum!complexity determination}
The gravitational complexity determines the structural type of
$\mathcal{H}(T)$:
\begin{center}
\small
\renewcommand{\arraystretch}{1.25}
\begin{tabular}{c@{\quad}l@{\quad}l}
\toprule
\emph{Class} & \emph{Laplace kernel $r^L(z)$} &
\emph{Shadow connection} \\
\midrule
$\mathbf{G}$ &
  $r^L(z) = \kappa/z^2$. &
  $\nabla^{\mathrm{hol}} = d - \kappa\,\omega_g$ \\
$\mathbf{L}$ &
  $r^L(z) = \kappa/z^2 + \mathfrak{C}_0/z^3$. &
  $\nabla^{\mathrm{hol}} = d - \kappa\,\omega_g -
  \mathfrak{C}$ \\
$\mathbf{C}$ &
  $r^L(z) = \kappa/z^2 + \mathfrak{Q}_0/z^4$. &
  $\nabla^{\mathrm{hol}} = d - \kappa\,\omega_g -
  \mathfrak{Q}$ \\
$\mathbf{M}$ &
  $r^L(z) = \sum_{r\geq 2} S_r/z^r$. &
  $\nabla^{\mathrm{hol}} = d -
  \sum_{r\geq 2} \mathrm{Sh}_r$ \\
\bottomrule
\end{tabular}
\end{center}
\end{theorem}

\begin{proof}
The binary shadow $r^L(z)$ (the Laplace kernel of the
$\lambda$-bracket) receives an arity-$r$ contribution as a
$z^{-r}$ pole; the collision residue
$r^{\mathrm{coll}}(z) = \Res^{\mathrm{coll}}_{0,2}(\Theta_\cA)$
has pole orders one lower.
Stabilization at $r_{\max}$ truncates the Laurent expansion.
\end{proof}


\begin{theorem}[Boundary-holographic complexity classification]
\label{thm:thqg-boundary-holographic-complexity}
\ClaimStatusProvedHere
\index{boundary-holographic complexity|textbf}
\index{shadow depth!boundary-holographic interpretation}
\index{gravitational partition function!complexity classification}
Let $\cT$ be a $3$d HT theory on $\C \times \R$ whose boundary
chiral algebra $\cA$ is modular Koszul.
The bulk is topological \textup{(}$3$d Chern--Simons or its
higher-spin generalisation; no local propagating degrees of
freedom\textup{)}.  The perturbative gravitational partition
function
\begin{equation}\label{eq:boundary-holographic-Z}
  Z^{\mathrm{pert}}(\cT)
  \;=\;
  \exp\Bigl(\,\sum_{g \geq 1} \hbar^{2g}\,F_g(\cT)\Bigr),
  \qquad
  F_g(\cT)
  \;=\;
  \int_{\overline{\cM}_g}
  \operatorname{Tr}^{\circ g}\bigl(\Theta_\cA^{(g)}\bigr),
\end{equation}
is completely determined by the boundary data: the OPE of
$\cA$, its transferred $A_\infty$ structure, and the
bar-intrinsic MC element $\Theta_\cA$
\textup{(}Theorem~\textup{\ref*{V1-thm:mc2-bar-intrinsic}}\textup{)}.
The shadow depth $r_{\max}(\cA)$ classifies the nonlinear
complexity of this boundary-determined partition function.
On any primary line~$L$, the weighted shadow generating function
$H(t) = \sum_{r \geq 2} r\,\mathrm{Sh}_r\,t^r$ satisfies
the Riccati algebraicity relation
\textup{(}Theorem~\textup{\ref*{V1-thm:riccati-algebraicity}}\textup{)}
\begin{equation}\label{eq:boundary-holographic-riccati}
  H(t)^2 \;=\; t^4\,Q_L(t),
  \qquad
  Q_L(t) = 4\kappa^2 + 12\kappa\alpha\,t
  + (9\alpha^2 + 2\Delta)\,t^2,
\end{equation}
with critical discriminant $\Delta = 8\kappa S_4$.
The arity-$r$ shadow first contributes to $F_g$ at genus
$g_{\min}(r) = \lfloor r/2 \rfloor + 1$
\textup{(}Corollary~\textup{\ref{cor:shadow-visibility-genus}}\textup{)}.
The four complexity classes are:
\begin{enumerate}[label=\textup{(\roman*)}]
\item \emph{Gaussian} ($\mathbf{G}$, $r_{\max} = 2$;
  $\Delta = 0$, $\alpha = 0$).
  The boundary $A_\infty$ structure is strictly
  coassociative: $m_k = 0$ for $k \geq 3$.
  The partition function is one-loop exact at every genus:
  \[
    F_g(\cT) = \kappa(\cA)\cdot\lambda_g^{\mathrm{FP}},
    \qquad
    \lambda_g^{\mathrm{FP}} =
    \frac{2^{2g-1}-1}{2^{2g-1}}\cdot
    \frac{|B_{2g}|}{(2g)!},
  \]
  where $B_{2g}$ are Bernoulli numbers; no higher-arity
  corrections exist.  The complementarity potential
  $S_\cA(x) = \tfrac{1}{2}\kappa\,x^2$ is Gaussian.
  \textup{(}Heisenberg $\cH_k$, lattice VOA $V_\Lambda$,
  free fermion.\textup{)}

\item \emph{Lie} ($\mathbf{L}$, $r_{\max} = 3$;
  $\Delta = 0$, $\alpha \neq 0$).
  The boundary OPE produces a cubic shadow
  $\mathfrak{C} = \kappa(x,[y,z]) \neq 0$;
  the quartic obstruction $o_4 =
  \tfrac{1}{2}\{\mathfrak{C},\mathfrak{C}\}_H$ vanishes
  by the Jacobi identity.
  The complementarity potential is cubic:
  $S_\cA(x) = \tfrac{1}{2}\kappa\,x^2 +
  \tfrac{1}{6}\mathfrak{C}_0\,x^3$.
  New corrections first appear at genus~$2$
  \textup{(}from $\mathrm{Sh}_3$\textup{)} and no
  higher-arity shadows contribute.
  \textup{(}Affine $V_k(\fg)$ at generic level.\textup{)}

\item \emph{Contact} ($\mathbf{C}$, $r_{\max} = 4$;
  stratum separation\textup{)}.
  The boundary $A_\infty$ structure has
  $\mathfrak{C} = 0$ and a quartic contact invariant
  $\mathfrak{Q} \neq 0$.
  The quintic obstruction $o_5 = 0$ by rank-one abelian
  rigidity: the self-bracket
  $\mathrm{Def}_{\mathrm{cyc}}^{(q)} \otimes
  \mathrm{Def}_{\mathrm{cyc}}^{(q)} \to
  \mathrm{Def}_{\mathrm{cyc}}^{(2q)}$ has
  $\dim\mathrm{Def}_{\mathrm{cyc}}^{(2q)} = 0$.
  One quartic correction to the Gaussian, realized at
  genus~$3$.
  \textup{(}$\beta\gamma$ on the weight-changing line.\textup{)}

\item \emph{Mixed} ($\mathbf{M}$, $r_{\max} = \infty$;
  $\Delta \neq 0$).
  Both $\mathfrak{C} \neq 0$ and $\mathfrak{Q} \neq 0$;
  the cubic source permanently generates higher obstructions:
  $\{\mathfrak{C}, \mathrm{Sh}_r\}_H \neq 0$ for all
  $r \geq 3$ forces $o_{r+1} \neq 0$ by induction.
  The complementarity potential is non-polynomial but
  convergent.  The shadow coefficients satisfy
  $S_r = \Delta \cdot R_r$ for all $r \geq 4$, with
  $R_r \in \mathbb{Q}(\alpha, \kappa, S_4)$, and
  $|S_r| \sim A\,\rho^r\,r^{-5/2}$ with growth rate
  \[
    \rho(\cA)
    \;=\;
    \frac{\sqrt{9\alpha^2 + 2\Delta}}{2|\kappa|}
  \]
  \textup{(}Theorem~\textup{\ref*{V1-thm:shadow-radius}}\textup{)}.
  The partition function receives genuinely new corrections
  at every genus from every arity: two new shadow
  coefficients $(S_{2g-2}, S_{2g-1})$ enter at each genus~$g$.
  \textup{(}Virasoro $\mathrm{Vir}_c$, $\mathcal{W}_N$ for
  $N \geq 3$.\textup{)}
\end{enumerate}
The classification is exhaustive
\textup{(}Theorem~\textup{\ref*{V1-thm:single-line-dichotomy}}\textup{)}
and detected at genus~$2$: the binary invariants
$\sigma = [\mathfrak{C} \neq 0]$ and
$\pi = [\mathfrak{Q} \neq 0]$ biject onto
$\{\mathbf{G}, \mathbf{L}, \mathbf{C}, \mathbf{M}\}$.
The bulk has no independent dynamics; the entire perturbative
gravitational content resides in the boundary OPE and its
modular extension~$\Theta_\cA$.
\end{theorem}

\begin{proof}
The bulk theory is topological: $3$d Chern--Simons has no
local propagating degrees of freedom.  Perturbative partition
functions are entirely determined by the boundary chiral
algebra (Brown--Henneaux for Virasoro,
Theorem~\ref{thm:brown-henneaux};
the Costello--Gwilliam construction for the general case;
see~\S\ref{subsec:gravity-physics}).
Equation~\eqref{eq:boundary-holographic-Z} follows from the
bar-intrinsic MC element
(Theorem~\ref*{V1-thm:mc2-bar-intrinsic}): the genus-$g$
MC component $\Theta_\cA^{(g)}$ is a section of the
$g$-fold sewing of the Hodge-weighted bar complex, and the
trace $\operatorname{Tr}^{\circ g}$ extracts the scalar
genus-$g$ free energy.

The Riccati algebraicity~\eqref{eq:boundary-holographic-riccati}
is the restriction of the all-arity master equation to a primary
line
(Theorem~\ref*{V1-thm:riccati-algebraicity}): the nonlinear
recursion among shadow coefficients is equivalent to the
algebraic relation $H^2 = t^4 Q_L$, making the shadow
generating function algebraic of degree~$2$.  The discriminant
$\Delta = 8\kappa S_4$ controls the analytic type: $\Delta = 0$
gives $\sqrt{Q_L}$ polynomial \textup{(}finite tower\textup{)};
$\Delta \neq 0$ gives $\sqrt{Q_L}$ irrational
\textup{(}infinite tower\textup{)}.

The shadow visibility genus
$g_{\min}(S_r) = \lfloor r/2 \rfloor + 1$ follows from the
self-loop parity vanishing
(Proposition~\ref*{V1-prop:self-loop-vanishing}): single-vertex
graphs $(0, 2k)$ with $k$ self-loops contribute zero for
$k \geq 2$ by odd-dimensional parity obstruction; the first
nonvanishing contribution of~$S_r$ comes from a two-vertex
graph at genus $\lfloor r/2 \rfloor + 1$.

Parts~(i)--(iv) then follow from the four-class classification
and the stabilization theorem above, together with the shadow
archetype classification
(Theorem~\ref*{V1-thm:shadow-archetype-classification}).
For class~$\mathbf{G}$, only the quadratic shadow~$\kappa$
contributes, giving
$F_g = \kappa \cdot \lambda_g^{\mathrm{FP}}$
(the scalar genus tower, Theorem~D).
For class~$\mathbf{M}$, the non-terminating tower contributes
corrections at every genus, with convergent generating function
by the shadow radius estimate
(Theorem~\ref*{V1-thm:shadow-radius}).
Exhaustiveness and genus-$2$ detection follow from the
single-line dichotomy
(Theorem~\ref*{V1-thm:single-line-dichotomy})
and the genus-$2$ binary diagnostic above: the map
$(\sigma, \pi) \to \{\mathbf{G}, \mathbf{L},
\mathbf{C}, \mathbf{M}\}$ is a bijection.
\end{proof}


% ======================================================================
%  §20.  THE $W_N$ LIMIT
% ======================================================================

\subsection{The $\mathcal{W}_\infty$ limit}
% label removed: subsec:thqg-winfty-limit
\index{W-infinity@$\mathcal{W}_\infty$!complexity limit}

\begin{proposition}[$\mathcal{W}_N$ complexity stabilises]
% label removed: prop:thqg-wn-stabilization
\ClaimStatusProvedHere
For all $N \geq 2$: $\mathcal{W}_N^k$ is class~$\mathbf{M}$
with $r_{\max} = \infty$.  The coefficients
$\mathrm{Sh}_r(\mathcal{W}_N^k)$ stabilise in $N$ for fixed~$r$:
$\mathrm{Sh}_r(\mathcal{W}_N) = \mathrm{Sh}_r(\mathcal{W}_\infty)$
for $N \geq r - 1$.
\end{proposition}

\begin{proof}
Each $\mathcal{W}_N$ for $N \geq 2$ contains Virasoro,
so $\mathfrak{C} \neq 0$ and $\mathfrak{Q} \neq 0$.
The infinite tower follows from
Lemma~\ref{lem:thqg-cubic-source}.
Stabilisation: the arity-$r$ shadow involves generators of
conformal weight $\leq r$; for $N > r-1$, these coincide
in $\mathcal{W}_N$ and $\mathcal{W}_{N+1}$.
\end{proof}


% ======================================================================
%  §21.  THE STANDARD LANDSCAPE
% ======================================================================

\subsection{The standard gravitational landscape}
% label removed: subsec:thqg-grav-landscape
\index{gravitational complexity!landscape}

\begin{theorem}[Gravitational complexity census]
% label removed: thm:thqg-grav-landscape
\ClaimStatusProvedHere
\index{gravitational complexity!complete census}
\begin{center}
\small
\renewcommand{\arraystretch}{1.3}
\begin{tabular}{l@{\quad}c@{\quad}c@{\quad}l@{\quad}l}
\toprule
\emph{Family} & \emph{Class} & $r_{\max}$ &
  \emph{Potential} &
  \emph{Bulk gravity} \\
\midrule
Heisenberg $\cH_k$ &
  $\mathbf{G}$ & $2$ & Quadratic & Free \\
Free fermion &
  $\mathbf{G}$ & $2$ & Quadratic & Free \\
Lattice $V_\Lambda$ &
  $\mathbf{G}$ & $2$ & Quadratic & Free \\[3pt]
$V_k(\mathfrak{sl}_2)$ &
  $\mathbf{L}$ & $3$ & Cubic & Semistrict HS \\
$V_k(\mathfrak{sl}_N)$ &
  $\mathbf{L}$ & $3$ & Cubic & Semistrict HS \\
$V_k(\fg)$ (simple) &
  $\mathbf{L}$ & $3$ & Cubic & Semistrict HS \\[3pt]
$\beta\gamma$ &
  $\mathbf{C}$ & $4$ & Quartic & Contact gravity \\[3pt]
$\mathrm{Vir}_c$ &
  $\mathbf{M}$ & $\infty$ & Non-poly & Full $L_\infty$ \\
$\mathcal{W}_3^k$ &
  $\mathbf{M}$ & $\infty$ & Non-poly & Full $L_\infty$ \\
$\mathcal{W}_N^k$ &
  $\mathbf{M}$ & $\infty$ & Non-poly & Full $L_\infty$ \\
\bottomrule
\end{tabular}
\end{center}
All families are Koszul; the complexity classification refines
within the Koszul locus.
\end{theorem}


% ======================================================================
%  §22.  THE MAIN THEOREM
% ======================================================================

\subsection{Summary: the gravitational complexity theorem}
% label removed: subsec:thqg-complexity-summary

\begin{maintheorem}[Gravitational complexity; result (G2)]
% label removed: thm:thqg-g2-main
\ClaimStatusProvedHere
\index{gravitational complexity!main theorem}
Let $\cA$ be a chirally Koszul, cyclically admissible chiral
algebra satisfying the HS-sewing condition.
The shadow depth $r_{\max}(\cA)$ determines the minimal
homotopy-algebraic complexity of the bulk gravitational theory.
The classification into exactly four classes is forced by the
obstruction tower:
\begin{enumerate}[label=\textup{(\roman*)}]
\item \textbf{Gaussian} ($r_{\max} = 2$):
  Free bulk theory.  All obstructions vanish.
  Quadratic potential.  Archetype: Heisenberg.

\item \textbf{Lie} ($r_{\max} = 3$):
  Semistrict higher-spin gravity with cubic vertices.
  Jacobi identity kills all higher arities.
  Archetype: affine $V_k(\fg)$.

\item \textbf{Contact} ($r_{\max} = 4$):
  Quartic contact vertices, no cubic.
  Cubic gauge triviality + rank-one rigidity terminate the tower.
  Archetype: $\beta\gamma$.

\item \textbf{Mixed} ($r_{\max} = \infty$):
  Full $L_\infty$ with vertices at all arities.
  Cubic source permanently generates higher obstructions.
  Non-polynomial potential.  Archetype: Virasoro.
\end{enumerate}
The gravitational complexity is a quasi-isomorphism invariant,
a Koszul duality invariant, and generically stable under
parameter deformation.
\end{maintheorem}


% ======================================================================
%  §23.  ENVELOPE-SHADOW FUNCTOR
% ======================================================================

\subsection{Gravitational complexity as a functor}
% label removed: subsec:thqg-envelope-shadow
\index{envelope-shadow functor!gravitational}

\begin{proposition}[Functoriality of complexity]
% label removed: prop:thqg-complexity-functor
\ClaimStatusProvedHere
\index{envelope-shadow functor!complexity functor}
For chiral algebras arising as $\cA = \Uvert(R)$
\textup{(}Nishinaka~\cite{Nish26}\textup{)}:
\begin{equation}% label removed: eq:thqg-chienv
\chienv(R)
\;:=\;
\rho_{\mathrm{grav}}(\Uvert(R))
\end{equation}
is a functor on Lie conformal algebras.  Three structural
properties hold:
\begin{enumerate}[label=\textup{(\roman*)}]
\item \emph{Finite-jet rigidity}
  \textup{(}Proposition~\ref{prop:finite-jet-rigidity}\textup{)}:
  the arity-$r$ shadow depends only on the $(r{+}1)$-jet
  $\Uvert(R)/F_{r+1}$.
\item \emph{Polynomial level dependence}
  \textup{(}Proposition~\ref{prop:polynomial-level-dependence}\textup{)}:
  the arity-$r$ shadow is polynomial of degree $\leq r{-}1$ in
  a central-extension parameter.
\item \emph{Gaussian collapse}
  \textup{(}Proposition~\ref{prop:gaussian-collapse-abelian}\textup{)}:
  $\chienv(R) = 2$ if and only if $R$ is abelian; Heisenberg is
  the universal fixed point.
\end{enumerate}
\end{proposition}

\begin{proof}
Functoriality: a morphism $R \to R'$ induces
$\Uvert(R) \to \Uvert(R')$, hence
$\Theta_{\Uvert(R)} \mapsto \Theta_{\Uvert(R')}$.
Shadow projections are compatible, giving
$\chienv(R) \leq \chienv(R')$ when $R \hookrightarrow R'$.
Part~(i): the obstruction class $o_{r+1}$ depends on
the structure constants of the $\lambda$-bracket through
degree $r+1$ (the graph-sum formula involves at most $r+1$
vertices), hence on the $(r{+}1)$-jet.
Part~(ii): the structure constants of $\Uvert(R)$ are polynomial
in the central extension $k$ (by the OPE expansion), so
$o_{r+1}$ is polynomial in $k$ of degree at most $r-1$
(each vertex contributes at most degree one).
Part~(iii): if $R$ is abelian, all iterated brackets vanish
and $d_\infty = 2$; conversely, $d_\infty = 2$ forces
all transferred ternary operations to vanish, which implies
$a_{(0)}b = 0$ for all generators.
\end{proof}

\begin{remark}[DS reduction and complexity]
% label removed: rem:thqg-ds-complexity
\index{Drinfeld--Sokolov reduction!complexity}
For the principal Drinfeld--Sokolov reduction
$\mathcal{W}_k(\fg) = H^0_{\mathrm{DS}}(V_k(\fg))$:
the reduction preserves Koszulness but may change the
complexity class.  Affine $V_k(\mathfrak{sl}_2)$ is
class~$\mathbf{L}$ ($r_{\max} = 3$), while its DS reduction
$\mathrm{Vir}_c$ is class~$\mathbf{M}$ ($r_{\max} = \infty$).
The reduction $V_k(\mathfrak{sl}_N) \to \mathcal{W}_N^k$
similarly jumps from $\mathbf{L}$ to $\mathbf{M}$ for
all $N \geq 2$.  The mechanism: DS reduction introduces
non-abelian composite fields (Sugawara $T$, higher-spin
$W_s$) whose iterated OPE generates operations at all
arities.
\end{remark}

\begin{remark}[Chain-level versus cohomology-level structure for $\mathcal{W}_N$]
\label{rem:ds-chain-vs-cohomology-wn}
\index{Drinfeld--Sokolov reduction!chain vs cohomology}
\index{W-algebra!chain-level structure}
The $\Ainf$ operations and the Yangian structure for
$\mathcal{W}_N$ live on the BRST \emph{cohomology}
$H^\bullet(V_k(\mathfrak{sl}_N) \otimes \mathcal{C}_f,\,
Q_{\mathrm{DS}})$, not on the chain-level BRST complex.  The
passage from chains to cohomology is the homotopy transfer
theorem (Kadeishvili--Merkulov), and it discards
information.  Three phenomena deserve attention.

\begin{enumerate}[label=\textup{(\roman*)}]
\item \emph{The chain-level complex is strictly richer.}
  For every $N \geq 2$, the BRST complex
  $\bigl(V_k(\mathfrak{sl}_N) \otimes \mathcal{C}_f,\,
  Q_{\mathrm{DS}}\bigr)$ is a dg vertex algebra whose
  generators include the affine currents $J^a(z)$ and the
  $bc$-ghost system $(\beta_\alpha, \gamma_\alpha)$ for each
  positive root~$\alpha$.  The cohomology
  $\mathcal{W}_N = H^0(Q_{\mathrm{DS}})$ has generators
  $W_2 = T,\, W_3,\, \ldots,\, W_N$ of conformal weights
  $2, 3, \ldots, N$.  The composite fields appearing in the
  $\mathcal{W}_N$ OPE (in particular the quasi-primary
  $\Lambda = {:\!TT\!:} - \tfrac{3}{10}\partial^2 T$ at
  weight~$4$ for $N = 3$) are \emph{not} among the
  generators of the chain-level complex; they arise only
  after taking $Q_{\mathrm{DS}}$-cohomology and forming
  normal-ordered products.  The chain-level differential
  $Q_{\mathrm{DS}}$ sees the individual $J^a$ and ghost
  contributions to~$\Lambda$ separately; the cohomological
  algebra sees only the composite.

\item \emph{Higher transferred operations detect the loss.}
  The chain-level $\Ainf$ structure
  $\{m_k^{\mathrm{chain}}\}$ on the BRST complex is strict
  (coming from the operadic action of~$\SCchtop$ on the
  gauge theory).  The transferred operations
  $\{m_k^H\}$ on $\mathcal{W}_N$ are computed by
  summing over trees with internal edges labelled by the
  BRST contracting homotopy~$h$.  For $N = 2$ (Virasoro),
  $h$ threads through the single ghost pair
  $(\beta, \gamma)$ and the Sugawara construction yields the
  wheel-diagram formula for $m_k^H(T, \ldots, T)$ at all
  arities~$k \geq 3$.  For $N = 3$, the homotopy~$h$ routes
  through $\binom{3}{2} = 3$ ghost pairs; the transferred
  $m_3^H(W, W, W)$ involves compositions
  $m_2^{\mathrm{chain}}(h \circ m_2^{\mathrm{chain}}, -)$
  whose intermediate states pass through ghost-number sectors
  invisible to~$\mathcal{W}_3$.  In particular, the
  $c$-dependent coefficient
  $\tfrac{32}{5c+22}$ multiplying~$\Lambda$ in the
  $\{W_\lambda W\}$ bracket is the \emph{residue} of a
  chain-level computation that also determines the
  relative coefficient of non-cohomological ghost terms;
  these ghost terms are projected away on cohomology but
  contribute to the transferred $m_4^H, m_5^H, \ldots\,$
  through the homotopy.

\item \emph{General $N$: the ghost sector grows
  quadratically.}
  For $\mathfrak{sl}_N$, the BRST complex has
  $\tfrac{1}{2}N(N-1)$ ghost pairs (one per positive root).
  The chain-level state space at conformal weight~$h$
  grows as
  $\dim\bigl(V_k(\mathfrak{sl}_N) \otimes
  \mathcal{C}_f\bigr)_h \sim p_N(h)$,
  where $p_N$ is a partition function with $N^2 - 1$
  free-field contributions.  The cohomology
  $\dim(\mathcal{W}_N)_h$ is governed by the $N - 1$
  generators of weights $2, \ldots, N$, hence is
  strictly smaller for $h \geq N + 1$.
  The homotopy transfer discards
  $\dim(\mathrm{chain})_h -
  \dim(\mathrm{cohomology})_h$
  directions at each weight; these discarded directions
  are the ``reservoir'' from which the transferred
  $m_k^H$ draws its non-trivial values.  Concretely:
  the transferred $m_k^H$ on $\mathcal{W}_N$ depends on
  the $Q_{\mathrm{DS}}$-acyclic summand through the
  contracting homotopy, and this dependence is sensitive
  to the choice of homotopy retract (different choices
  yield gauge-equivalent $\Ainf$ structures).
\end{enumerate}

For non-principal reductions
$\mathcal{W}_k(\fg, f) = H^0(Q_{\mathrm{DS}}^f)$
with $f$ not the principal nilpotent, the chain-level
complex is richer still: the ghost system
$\mathcal{C}_f$ is determined by the $\mathrm{ad}(x)$-eigenspace
decomposition ($x$ the neutral element of the
$\mathfrak{sl}_2$-triple containing~$f$), and
non-principal choices leave more affine generators in
cohomology while simultaneously enlarging the ghost
sector.  The subregular reduction for $\mathfrak{sl}_3$
produces a $\mathcal{W}$-algebra with generators at
weights $2$ and $3$ (as for $\mathcal{W}_3$) but with
different OPE coefficients and a different contracting
homotopy; the hook-type reductions in type~A produce
generators whose weights are determined by the
corresponding partition.  In every case, the
$\Ainf$ operations on the reduced algebra are shadows
of the chain-level BRST complex, and the full
chain-level data, accessible through the bar complex
$B(V_k(\fg) \otimes \mathcal{C}_f)$ before taking
cohomology, encodes strictly more information.

The open question is whether the chain-level BRST
data can be recovered from the cohomology-level
$\Ainf$ structure.  For principal reductions of simply
laced algebras, the answer is affirmative on the Koszul
locus: bar-cobar inversion (Theorem~A) reconstructs
the chain-level algebra from its transferred $\Ainf$
structure up to quasi-isomorphism.  For non-principal
reductions, the reconstruction requires that the
Koszul property survives the reduction
(Theorem~\ref*{V1-thm:ds-koszul-obstruction}), which is
proved for the hook-type corridor in type~A and open in
general.
\end{remark}


% ======================================================================
%  §24.  CONVERGENCE OF THE VIRASORO POTENTIAL
% ======================================================================

\subsection{Convergence analysis}
% label removed: subsec:thqg-convergence
\index{Virasoro!potential convergence}

\begin{proposition}[Coefficient asymptotics]
% label removed: prop:thqg-coefficient-asymptotics
\ClaimStatusProvedHere
\index{Virasoro!shadow tower asymptotics}
For large $c > 0$ and $r \geq 4$, the Virasoro shadow
coefficients $S_r$ (where $\mathrm{Sh}_r = S_r\,x^r$)
satisfy
\begin{equation}% label removed: eq:thqg-coefficient-bound
|S_r|
\;\leq\;
C(c)\cdot r!\cdot c^{-(r-3)}\cdot(5c+22)^{-\lfloor(r-2)/2\rfloor}
\end{equation}
for a constant $C(c)$ independent of~$r$.  In particular:
\begin{enumerate}[label=\textup{(\roman*)}]
\item For fixed $c > 0$, $|S_r| \to 0$ faster than any
  exponential as $r \to \infty$.
\item The complementarity potential
  $S_{\mathrm{Vir}}(x)$ converges absolutely in
  $|x| < R(c)$ with radius
  $R(c) \sim c^{1/2}$ for large~$c$.
\item The $c \to \infty$ limit of the normalised potential
  $c^{-1}\,S_{\mathrm{Vir}}(x\,c^{1/2})$ is the Gaussian
  $\frac{1}{4}\,y^2$: the large-$c$ limit is class~$\mathbf{G}$.
\end{enumerate}
\end{proposition}

\begin{proof}
Part~(i): The denominator pattern
(Proposition~\ref{prop:thqg-virasoro-structure}(i)) gives
$|S_r| \leq C_0\,r^{C_1}/[c^{r-3}(5c+22)^{\lfloor(r-2)/2\rfloor}]$
for polynomial numerator growth.  The double-exponential
decay in $r$ from the denominator dominates.
Part~(ii): The ratio test on $|S_{r+1}/S_r|$ gives
$\limsup_{r\to\infty} |S_{r+1}/S_r|^{1/r} \leq c^{-1}$,
so the radius of convergence is at least~$c$.
Part~(iii): Rescaling $x = y\,c^{-1/2}$ in the potential
and taking $c \to \infty$: the quadratic term
$\frac{c}{4}(y/\sqrt{c})^2 = \frac{1}{4}y^2$ survives;
the cubic term $\frac{1}{3}(y/\sqrt{c})^3 = O(c^{-3/2})$
vanishes; all higher terms decay faster.
\end{proof}

\begin{remark}[Singular loci]
% label removed: rem:thqg-singular-loci
\index{Virasoro!singular loci}
The potential $S_{\mathrm{Vir}}$ has two singular loci:
\begin{enumerate}[label=\textup{(\alph*)}]
\item $c = 0$: the Hessian $H = c/2$ degenerates, the
  propagator $P = 2/c$ diverges, and all shadow coefficients
  for $r \geq 4$ blow up.  Physically: the free-field
  limit where conformal symmetry is trivialised.
\item $c = -22/5$: the quartic contact coefficient
  $Q_0 = 10/[c(5c+22)]$ diverges.  This is the minimal
  model $M(2,5)$ central charge.  The collision of Shapovalov
  eigenvalues at level~$2$ produces the singularity.
\end{enumerate}
At both loci, the shadow obstruction tower requires analytic continuation,
not direct evaluation.
\end{remark}


% ======================================================================
%  §25.  DEFORMATION STABILITY
% ======================================================================

\subsection{Deformation stability of the complexity class}
% label removed: subsec:thqg-deformation-stability
\index{shadow depth!deformation stability}

\begin{proposition}[Generic constancy]
% label removed: prop:thqg-generic-constancy
\ClaimStatusProvedHere
\index{shadow depth!generic constancy}
Let $\{\cA_t\}_{t \in \C}$ be a continuous family of modular
Koszul chiral algebras.  The set
$\{t : r_{\max}(\cA_t) \neq r_{\max}(\cA_0)\}$ is discrete
in~$\C$.  At critical values, $r_{\max}$ can only decrease
\textup{(}obstruction classes specialise to zero\textup{)}.
\end{proposition}

\begin{proof}
The obstruction class $o_{r+1}(\cA_t)$ is a polynomial in
the structure constants of $\cA_t$ (graph-sum formula).  As
a function of $t$, it is algebraic.  The zero locus of an
algebraic function is discrete unless identically zero.
Vanishing at a special value is a closed condition, so
$r_{\max}$ can only decrease.
\end{proof}

\begin{example}[Class jumps in the Virasoro family]
% label removed: ex:thqg-virasoro-class-jumps
\index{Virasoro!class jumps}
The Virasoro algebra $\mathrm{Vir}_c$ is class~$\mathbf{M}$
for all $c \neq 0$, with no class jumps at finite $c$.
The quartic contact $Q_0 = 10/[c(5c+22)]$ is nonzero for
$c \neq 0, -22/5$, and the cubic $\mathfrak{C} = 2x^3$ is
independent of $c$.  At $c = 0$: the Hessian degenerates;
the shadow depth is undefined (the cyclic pairing fails to
be nondegenerate).  At $c = -22/5$: $Q_0$ diverges, signalling
a change of the effective field space (the minimal model
$M(2,5)$ has a null vector that modifies the cyclic slice).
\end{example}

\begin{example}[Class jumps in the affine family]
% label removed: ex:thqg-affine-class-jumps
\index{affine Kac--Moody!class jumps}
The affine algebra $V_k(\mathfrak{sl}_2)$ is class~$\mathbf{L}$
for all non-critical levels $k \neq -2$.  The cubic shadow
$\mathfrak{C} = \kappa \cdot f_{\mathfrak{sl}_2}\cdot x^3$
vanishes at $k = -h^\vee = -2$ (critical level, where
$\kappa = 0$).  At $k = -2$: the shadow depth jumps from
$3$ to $2$ (class $\mathbf{L} \to \mathbf{G}$), consistent
with the Feigin--Frenkel construction at the critical level.
\end{example}


% ======================================================================
%  §26.  THE MC EQUATION AS GRAVITY ACTION
% ======================================================================

\subsection{The MC equation as a gravitational action principle}
% label removed: subsec:thqg-mc-action
\index{Maurer--Cartan equation!as gravitational action}

\begin{theorem}[MC equation = Euler--Lagrange equation]
% label removed: thm:thqg-mc-euler-lagrange
\ClaimStatusProvedHere
\index{complementarity potential!as gravity action}
The MC equation $D_\cA\,\Theta_\cA +
\frac{1}{2}[\Theta_\cA, \Theta_\cA] = 0$,
evaluated on the one-generator primary line, is the
Euler--Lagrange equation for the complementarity potential
$S_\cA(x)$:
\begin{equation}% label removed: eq:thqg-euler-lagrange
\frac{\partial S_\cA}{\partial x}
\;=\;
\kappa\,x + \sum_{r=3}^{r_{\max}}
\mathrm{Sh}_r(\cA)\,x^{r-1}
\;=\; 0
\end{equation}
at the critical point $x = 0$.  The second variation is
$S''_\cA(0) = \kappa(\cA) = H_\cA$.  The genus-$g$
partition function arises from the formal saddle-point
expansion $Z_g(\cA) \sim \int e^{S_\cA/\hbar}$.
\end{theorem}

\begin{proof}
The all-arity master equation
$\nabla_H(\mathrm{Sh}_r) + \mathfrak{o}^{(r)} = 0$,
restricted to the primary line, is the recursion for the
Taylor coefficients of $S_\cA$.  The MC equation for
$\Theta_\cA$ is equivalent to
$\frac{1}{2}\{S, S\}_{H_\cA} + \hbar\,\Delta\,S = 0$
(the BV classical master equation).  At genus~$0$ ($\hbar = 0$):
$\frac{1}{2}\{S, S\}_H = 0$, which is the Euler--Lagrange
condition for the BV action.
\end{proof}

\begin{remark}[The four potentials as four gravity theories]
% label removed: rem:thqg-four-gravity-theories
\index{gravitational action!four types}
\begin{center}
\small
\renewcommand{\arraystretch}{1.25}
\begin{tabular}{cll}
\toprule
\emph{Class} & $S_\cA(x)$ &
\emph{Gravity analogue} \\
\midrule
$\mathbf{G}$ &
  $\frac{\kappa}{2}\,x^2$ &
  Free massless field (Gaussian integral) \\
$\mathbf{L}$ &
  $\frac{\kappa}{2}\,x^2 + \frac{1}{3!}\mathfrak{C}_0\,x^3$ &
  Chern--Simons gravity ($A\wedge dA + \frac{2}{3}A^3$) \\
$\mathbf{C}$ &
  $\frac{1}{4!}\mathfrak{Q}_0\,x^4$ &
  $\phi^4$ contact theory \\
$\mathbf{M}$ &
  $\sum_{r=2}^\infty \frac{1}{r}S_r\,x^r$ &
  Full quantum gravity (non-polynomial EFT) \\
\bottomrule
\end{tabular}
\end{center}
The classification of gravitational effective actions
by their polynomial degree mirrors the standard EFT
truncation hierarchy: renormalisable ($r \leq 4$) vs.\
non-renormalisable ($r = \infty$).  Class~$\mathbf{M}$
is the gravitational analogue of a non-renormalisable
theory requiring infinitely many counterterms in the
Wilsonian sense, but here the full tower is determined
by the finite OPE data of the boundary algebra, so no
independent counterterms arise.
\end{remark}


% ======================================================================
%  §27.  THE HOLOGRAPHIC SPECTRAL SEQUENCE
% ======================================================================

\subsection{The holographic spectral sequence}
% label removed: subsec:thqg-holographic-ss
\index{holographic spectral sequence|textbf}

The collision filtration on the convolution algebra
$\gAmod$ induces a spectral sequence that separates the
contributions of the four complexity classes.

\begin{construction}[Collision filtration]
% label removed: constr:thqg-collision-filtration
\index{collision filtration|textbf}
The arity filtration
$F^r\gAmod := \bigoplus_{s \geq r}(\gAmod)_s$
defines a complete, exhaustive, separated filtration
of the dg Lie algebra.  The associated graded is
\begin{equation}% label removed: eq:thqg-graded-convolution
\operatorname{gr}^r_F(\gAmod)
\;=\;
(\gAmod)_r
\;\cong\;
\Defcyc(\cA)_r \widehat\otimes \Gmod_r,
\end{equation}
the arity-$r$ piece of the convolution algebra.
\end{construction}

\begin{theorem}[Holographic spectral sequence]
% label removed: thm:thqg-holographic-ss
\ClaimStatusProvedHere
\index{spectral sequence!holographic}
The collision filtration induces a convergent spectral sequence
\begin{equation}% label removed: eq:thqg-ss
E_1^{r,q}
\;=\;
H^{r+q}(\operatorname{gr}^r_F(\gAmod),\, d_2)
\;\Longrightarrow\;
H^{r+q}(\gAmod).
\end{equation}
The differential $d_1\colon E_1^{r,*} \to E_1^{r+1,*}$ is
the obstruction map $o_{r+1}$.  The $E_2$ page is
\begin{equation}% label removed: eq:thqg-e2
E_2^{r,q}
\;=\;
\ker(o_{r+1})/\mathrm{im}(o_r)
\;\text{ at the $(r,q)$ position.}
\end{equation}
\end{theorem}

\begin{proof}
Standard spectral sequence of a filtered dg Lie algebra
(Theorem~\ref{thm:spectral-sequence-filtered-dg}).  The
$E_1$ page computes the cohomology of each graded piece
with respect to the twisted differential $d_2$.  The
$d_1$ differential is the boundary map in the long exact
sequence of the filtration, which is exactly the
obstruction class $o_{r+1}$ at the $(r, r+1)$ position.
Convergence follows from completeness and pronilpotence of
$F^2\gAmod$.
\end{proof}

\begin{proposition}[Collapse criteria by complexity class]
% label removed: prop:thqg-collapse-criteria
\ClaimStatusProvedHere
\index{spectral sequence!collapse criteria}
\begin{enumerate}[label=\textup{(\roman*)}]
\item \emph{Class~$\mathbf{G}$}: The spectral sequence
  collapses at $E_1$.  All differentials $d_r = 0$ for
  $r \geq 1$.  The MC element is purely scalar in arity:
  $\Theta_\cA=\eta\otimes\Gamma_\cA$ for some genus
  coefficient~$\Gamma_\cA\in\widehat{\Gmod}$.  On the proved
  scalar lane this further specializes to
  $\Theta_\cA = \kappa\cdot\eta\otimes\Lambda$.
\item \emph{Class~$\mathbf{L}$}: The spectral sequence
  collapses at $E_2$.  The differential $d_1\colon
  E_1^{2,*} \to E_1^{3,*}$ is nonzero (it is $o_3$);
  $d_1\colon E_1^{3,*} \to E_1^{4,*}$ is zero (it is $o_4 = 0$).
  All $d_r = 0$ for $r \geq 2$.
\item \emph{Class~$\mathbf{C}$}: The spectral sequence
  collapses at $E_2$.  The differential $d_1\colon
  E_1^{3,*} \to E_1^{4,*}$ is nonzero ($o_4 \neq 0$);
  $d_1\colon E_1^{4,*} \to E_1^{5,*}$ is zero ($o_5 = 0$).
\item \emph{Class~$\mathbf{M}$}: The spectral sequence does
  not collapse at any finite page.  There are nonzero
  differentials $d_1$ at infinitely many positions.  The
  abutment is computed by the full inverse limit
  $\Theta_\cA = \varprojlim_r \Theta_\cA^{\leq r}$.
\end{enumerate}
\end{proposition}

\begin{proof}
Parts~(i)--(iii): in each case, the collapse page is
determined by the highest arity at which $o_{r+1} \neq 0$.
For class~$\mathbf{G}$: all $o_r = 0$ for $r \geq 3$,
so $d_1 = 0$ everywhere, giving $E_1$-collapse.
For class~$\mathbf{L}$: $o_3 \neq 0$ but $o_r = 0$ for
$r \geq 4$; the single nonzero differential at arity~$3$
produces the cubic shadow as a nontrivial $E_2$-class,
then all higher differentials vanish.
For class~$\mathbf{C}$: the gap ($o_3 = 0$, $o_4 \neq 0$,
$o_5 = 0$) gives $d_1$ nonzero at position~$3 \to 4$ and
zero elsewhere, again yielding $E_2$-collapse.
Part~(iv): for class~$\mathbf{M}$, the infinite tower
(Theorem~\ref{thm:thqg-virasoro-infinite}) gives
$o_{r+1} \neq 0$ for infinitely many~$r$, so $d_1$ has
infinitely many nonzero entries and the spectral sequence
never degenerates.
\end{proof}


% ======================================================================
%  §28.  THE $E_1$ PAGE: EXPLICIT STRUCTURE
% ======================================================================

\subsection{The $E_1$ page for Virasoro}
% label removed: subsec:thqg-e1-virasoro
\index{holographic spectral sequence!$E_1$ for Virasoro}

\begin{proposition}[$E_1$ page for Virasoro]
% label removed: prop:thqg-e1-virasoro
\ClaimStatusProvedHere
For the Virasoro algebra on the one-generator primary line:
\begin{equation}% label removed: eq:thqg-e1-virasoro
E_1^{r,0}
\;\cong\;
\C\cdot x^r
\quad\text{for all } r \geq 2,
\end{equation}
with differential $d_1(S_r\,x^r) = o_{r+1}(\mathrm{Vir})
\cdot x^{r+1}$.  The $E_2$ page is
\begin{equation}% label removed: eq:thqg-e2-virasoro
E_2^{r,0} = 0
\quad\text{for all } r \geq 2.
\end{equation}
Every $E_1$-class is killed by a differential: the spectral
sequence converges to a point.  This is consistent with the
bar-intrinsic construction: $\Theta_{\mathrm{Vir}}$ exists
and is unique, so the abutment $H^*(\gAmod)$ carries no
residual ambiguity.
\end{proposition}

\begin{proof}
On the rank-one primary line,
$\operatorname{gr}^r_F(\gAmod) \cong \C\cdot x^r$ in
degree zero.  The twisted differential $d_2$ acts by
multiplication by $\kappa = c/2$ and has no kernel beyond
degree zero.  The $d_1$ differential
$d_1(x^r) = o_{r+1}\cdot x^{r+1}$ is nonzero for all
$r \geq 3$ (since $o_{r+1} \neq 0$ for all $r \geq 4$
by the infinite tower), so
$E_2^{r,0} = \ker(d_1^{r,0})/\mathrm{im}(d_1^{r-1,0}) = 0$
for $r \geq 3$: every class at arity~$r$ is the image of
the arity-$(r{-}1)$ class under $d_1$ (up to a nonzero
scalar), and is simultaneously in the kernel of $d_1$ at
arity~$r$ (since $o_{r+2} \neq 0$ means $d_1$ is injective).
\end{proof}


% ======================================================================
%  §29.  THE $E_1$ COLLAPSE FOR AFFINE AND $\beta\gamma$
% ======================================================================

\subsection{Spectral sequence collapse for finite-depth classes}
% label removed: subsec:thqg-ss-collapse

\begin{proposition}[Collapse for class~$\mathbf{L}$]
% label removed: prop:thqg-collapse-L
\ClaimStatusProvedHere
\index{holographic spectral sequence!collapse for affine}
For $\cA = V_k(\fg)$ at generic $k$:
the spectral sequence collapses at $E_2$ with
\begin{equation}% label removed: eq:thqg-collapse-L
E_2^{2,0} \cong \C,
\qquad
E_2^{3,0} \cong \C,
\qquad
E_2^{r,0} = 0 \;\text{for}\; r \geq 4.
\end{equation}
The two surviving classes are $\kappa$ and $\mathfrak{C}$.
\end{proposition}

\begin{proof}
$d_1\colon E_1^{3,0} \to E_1^{4,0}$ is given by $o_4 = 0$
(Jacobi), so $E_1^{3,0}$ survives to $E_2$.
$d_1\colon E_1^{r,0} \to E_1^{r+1,0}$ is zero for all
$r \geq 3$ (no higher obstructions).  The $E_2$ page has
exactly two nonzero entries at genus zero.
\end{proof}

\begin{proposition}[Collapse for class~$\mathbf{C}$]
% label removed: prop:thqg-collapse-C
\ClaimStatusProvedHere
\index{holographic spectral sequence!collapse for $\beta\gamma$}
For $\cA = \beta\gamma$ on the weight-changing line:
the spectral sequence collapses at $E_2$ with
\begin{equation}% label removed: eq:thqg-collapse-C
E_2^{2,0} \cong \C,
\qquad
E_2^{4,0} \cong \C,
\qquad
E_2^{r,0} = 0 \;\text{for}\; r \neq 2, 4.
\end{equation}
The surviving classes are $\kappa$ and $\mathfrak{Q}$.
The gap ($E_2^{3,0} = 0$) is the spectral-sequence
manifestation of the cubic gauge triviality.
\end{proposition}

\begin{proof}
$o_3 = 0$ (abelian OPE), so $d_1\colon E_1^{2,0} \to E_1^{3,0}$
is zero and $E_1^{3,0}$ is not killed; but $o_4 \neq 0$, so
$d_1\colon E_1^{3,0} \to E_1^{4,0}$ is nonzero, killing
$E_1^{3,0}$ at the $E_2$ page (it becomes the image of
$d_1$, which maps to the kernel of the next $d_1$).
More precisely: $E_1^{3,0} \cong \C$ with $d_1 = 0$ from
arity~$2$ and $d_1 \neq 0$ to arity~$4$.  Since $o_3 = 0$,
$\mathrm{im}(d_1^{2,0}) = 0$, but
$\ker(d_1^{3,0}) = 0$ since $o_4 \neq 0$
gives an injective map.  Hence $E_2^{3,0} = 0$.
At arity~$4$: $d_1^{4,0} = 0$ ($o_5 = 0$ by rank-one
rigidity), and $\mathrm{im}(d_1^{3,0}) \neq 0$; but the
quotient $\ker(d_1^{4,0})/\mathrm{im}(d_1^{3,0})$ is
one-dimensional (the quartic class $\mathfrak{Q}$ modulo
the image of the cubic).
\end{proof}


% ======================================================================
%  §30.  CURVE INDEPENDENCE
% ======================================================================

\subsection{Universality: curve independence of complexity}
% label removed: subsec:thqg-curve-independence
\index{shadow depth!curve independence}

\begin{proposition}[Curve independence]
% label removed: prop:thqg-curve-independence
\ClaimStatusProvedHere
\index{gravitational complexity!curve independence}
The gravitational complexity $\rho_{\mathrm{grav}}(\cA)$ is
independent of the curve~$X$ on which $\cA$ lives.  The
obstruction classes $o_r(\cA)$ depend only on the OPE
coefficients, which are curve-independent data.
\end{proposition}

\begin{proof}
By the genus-$0$ curve-independence criterion
(Proposition~\ref{prop:genus0-curve-independence}):
the genus-$0$ shadow algebra $\cA^{\mathrm{sh}}_{r,0}$
depends only on the $\lambda$-bracket and the cyclic
pairing.  These are intrinsic to the vertex algebra,
not to the curve.
\end{proof}


% ======================================================================
%  §31.  EXPLICIT SEXTIC COMPUTATION
% ======================================================================

\subsection{Detailed sextic computation for Virasoro}
% label removed: subsec:thqg-sextic-detail
\index{Virasoro!sextic shadow detailed}

\begin{computation}[The sextic obstruction in full detail]
% label removed: comp:thqg-sextic
\index{obstruction class!sextic Virasoro}
The arity-$6$ obstruction $\mathfrak{o}^{(6)}$ receives
contributions from all two-vertex graphs with arities
$(a,b)$ satisfying $a + b = 7$, $a,b \geq 2$.  The
contributing pairs and their values on the Virasoro
primary line with $P = 2/c$:

\smallskip
\noindent
$(2,5)$: $\{\mathrm{Sh}_2, \mathrm{Sh}_5\}_H
= cx\cdot(2/c)\cdot(-240/[c^2(5c+22)])x^4
= -480/[c^2(5c+22)]\,x^5 \cdot x$. \\
$(3,4)$: $\{\mathfrak{C}, \mathfrak{Q}\}_H
= 6x^2\cdot(2/c)\cdot 4Q_0 x^3
= 48Q_0/c\,x^5\cdot x$. Already computed as part of
$\mathfrak{o}^{(5)}$; the arity-$6$ component requires
different combinatorics. \\
$(3,5)$: $\{\mathfrak{C}, \mathrm{Sh}_5\}_H
= 6x^2\cdot(2/c)\cdot 5S_5 x^4
= 60S_5/c\,x^6$.
With $S_5 = -48/[c^2(5c+22)]$:
$= -2880/[c^3(5c+22)]\,x^6$. \\
$(4,4)$: $\{\mathrm{Sh}_4, \mathrm{Sh}_4\}_H
= 4Q_0 x^3\cdot(2/c)\cdot 4Q_0 x^3
= 32Q_0^2/c\,x^6$.
With $Q_0 = 10/[c(5c+22)]$:
$= 3200/[c^3(5c+22)^2]\,x^6$. \\
$(5,3)$: $\{\mathrm{Sh}_5, \mathfrak{C}\}_H
= 5S_5 x^4\cdot(2/c)\cdot 6x^2
= 60S_5/c\,x^6 = -2880/[c^3(5c+22)]\,x^6$. \\
$(5,2)$: $\{\mathrm{Sh}_5, \mathrm{Sh}_2\}_H$:
this contributes to the genus-loop sector, not the
bracket term.  The arity-$2$ piece $\mathrm{Sh}_2 = (c/2)x^2$
paired with arity-$5$ gives $5S_5 x^4\cdot(2/c)\cdot cx
= 10S_5\,x^5$, but this produces arity $5 + 2 - 1 = 6$
with different symmetry.

\smallskip\noindent
Collecting the three dominant contributions at arity~$6$:
\begin{align*}
\mathfrak{o}^{(6)}
&= \{\mathfrak{C}, \mathrm{Sh}_5\}_H
+ \{\mathrm{Sh}_4, \mathrm{Sh}_4\}_H
+ \{\mathrm{Sh}_5, \mathfrak{C}\}_H
\\
&= -\frac{2880}{c^3(5c+22)}\,x^6
+ \frac{3200}{c^3(5c+22)^2}\,x^6
- \frac{2880}{c^3(5c+22)}\,x^6
\\
&= \frac{-5760(5c+22) + 3200}{c^3(5c+22)^2}\,x^6
\\
&= \frac{-28800c - 126720 + 3200}{c^3(5c+22)^2}\,x^6
\\
&= \frac{-28800c - 123520}{c^3(5c+22)^2}\,x^6.
\end{align*}
Factoring the numerator:
$-28800c - 123520 = -160(180c + 772) = -640(45c + 193)$.
Hence
\begin{equation}% label removed: eq:thqg-o6-explicit
\mathfrak{o}^{(6)}
= \frac{-640(45c+193)}{c^3(5c+22)^2}\,x^6.
\end{equation}
Applying the master recursion~\eqref{eq:thqg-master-recursion}:
\begin{equation}% label removed: eq:thqg-sh6-explicit
\mathrm{Sh}_6
= -\frac{\mathfrak{o}^{(6)}}{2\cdot 6}
= \frac{640(45c+193)}{12\,c^3(5c+22)^2}\,x^6
= \frac{80(45c+193)}{3\,c^3(5c+22)^2}\,x^6.
\end{equation}
This confirms the entry in the table of
Theorem~\ref{thm:thqg-virasoro-tower-explicit}.

The numerator $45c + 193$ vanishes at $c = -193/45$.  At this
value, $\mathrm{Sh}_6 = 0$ but $\mathrm{Sh}_5 \neq 0$ and
$\mathrm{Sh}_7 \neq 0$ (since $15c + 61|_{c=-193/45} =
15(-193/45) + 61 = -193/3 + 61 = (-193 + 183)/3 = -10/3 \neq 0$).
The shadow obstruction tower has a ``zero-crossing'' at arity~$6$ but remains
generically infinite.
\end{computation}


% ======================================================================
%  §32.  OPEN PROBLEMS
% ======================================================================

\subsection{Open problems}
% label removed: subsec:thqg-open-problems
\index{gravitational complexity!open problems}

\begin{openproblem}[Gap question]
% label removed: open:thqg-gap
\index{shadow depth!gap question}
Does there exist a modular Koszul chiral algebra with
$r_{\max} \in \{5, 6, 7, \ldots\} \setminus \{\infty\}$?
All known examples have $r_{\max} \in \{2, 3, 4, \infty\}$.
The exhaustiveness proof of
Theorem~\ref{thm:thqg-four-class} uses properties specific
to the standard landscape (rank-one rigidity for
class~$\mathbf{C}$).  An intermediate class would require:
$o_3 = 0$, $o_4 \neq 0$, $o_5 \neq 0$, $o_6 = 0$: a
quartic-plus-quintic tower that terminates at arity~$5$.
No known algebraic mechanism produces this pattern.
\end{openproblem}

% ----------------------------------------------------------------------
%  Shadow depth as stratification depth of the Steinberg
% ----------------------------------------------------------------------

\subsection{Shadow depth as stratification depth of the Steinberg}
% label removed: subsec:shadow-steinberg-stratification
\index{Steinberg object!stratification depth}
\index{shadow depth!Steinberg interpretation}

The shadow depth classification acquires a clean geometric
interpretation through the Chriss--Ginzburg picture: shadow
depth \emph{is} the stratification depth of the Steinberg
object $\mathfrak{S}_b$ associated to the bar complex.

\begin{definition}[Stratification depth]
% label removed: def:stratification-depth
Let $\mathfrak{S}_b$ denote the Steinberg object associated
to a modular Koszul chiral algebra~$\cA$.  The
\emph{stratification depth} of $\mathfrak{S}_b$ is the
length of the orbit stratification: the maximal~$d$ such
that there exist strata
\[
  S_0 \subsetneq S_1 \subsetneq \cdots \subsetneq S_d
  = \mathfrak{S}_b
\]
with each successive quotient $S_i / S_{i-1}$ non-empty.
\end{definition}

\begin{proposition}[Shadow depth = stratification depth]
% label removed: prop:shadow-equals-stratification
\ClaimStatusHeuristic
The shadow depth $r_{\max}$ of a modular Koszul chiral
algebra equals the stratification depth of the associated
Steinberg object~$\mathfrak{S}_b$.  Concretely:
\begin{enumerate}[label=\textup{(\roman*)}]
\item \textbf{Class~$\mathbf{G}$}
  ($r_{\max} = 2$): $\mathfrak{S}_b$ has two strata,
  the smooth locus plus one singular stratum.
  Examples: Heisenberg, lattice VOA.
\item \textbf{Class~$\mathbf{L}$}
  ($r_{\max} = 3$): three strata.
  Example: affine $V_k(\fg)$, where the Lie bracket
  creates one additional singular stratum.
\item \textbf{Class~$\mathbf{C}$}
  ($r_{\max} = 4$): four strata
  (the gap at depth~$3$ from cubic gauge triviality).
  Example: $\beta\gamma$.
\item \textbf{Class~$\mathbf{M}$}
  ($r_{\max} = \infty$): infinitely many strata.
  Example: Virasoro, $\mathcal{W}_N$ for $N \geq 3$, where
  higher-spin poles produce unbounded stratification.
\end{enumerate}
\end{proposition}

\begin{proof}[Argument]
Each stratum of $\mathfrak{S}_b$ corresponds to a distinct
collision type in $\mathrm{FM}_k(\mathbb{C})$ at which the
bar differential acquires a new pole order.  The bar
differential $d_B$ is assembled from OPE residues on
$\mathrm{FM}_k(\mathbb{C})$; a stratum $S_i$ is
\emph{non-trivial} precisely when the corresponding
collision produces a pole of order $\leq i$ that does not
factor through lower strata.  The identification of shadow
operations $\mathrm{Sh}_r$ with these strata rests on
the following chain: $\mathrm{Sh}_r \neq 0$ detects a
non-vanishing residue at pole order $r$ in the OPE, which
in turn produces a non-empty collision stratum at depth~$r$
in $\mathrm{FM}_k(\mathbb{C})$.  The converse (that every
non-empty stratum is detected by some $\mathrm{Sh}_r$)
follows from the completeness of the residue calculus on FM
strata.  We verify this in all three standard classes
(items~(i)--(iii) above); the general statement requires the
FM residue completeness axiom (H3).
\end{proof}

\begin{proposition}[Geometric translation of the gap question]
% label removed: prop:gap-steinberg
\ClaimStatusHeuristic
The gap question \textup{(}Open
Problem~\ref{open:thqg-gap}\textup{)} translates to: does
there exist a Steinberg object $\mathfrak{S}_b$ with
finitely many strata but more than four?

This is constrained by the representation theory of the
Fulton--MacPherson operad: the collision types in
$\mathrm{FM}_k(\mathbb{C})$ are classified by nested
parenthesizations, and the pole structure of the OPE
determines which collision types are non-trivial.  For a
polynomial OPE of degree~$d$ (pole of order~$d$), the
number of distinct collision strata grows at most linearly
in~$d$.  Finite shadow depth $> 4$ therefore requires a
very specific OPE pole pattern: poles of exactly orders
$1$ through $r_{\max}$ with all higher poles vanishing.
This is an extremely rigid condition on the chiral algebra.
\end{proposition}

\begin{proof}
The collision strata of $\mathrm{FM}_k(\mathbb{C})$ are
indexed by trees; each internal edge carries a pole order
from the OPE.  If the OPE has poles of orders
$1, \ldots, d$ and no higher, then the maximal
stratification depth is at most~$d$: each new pole order
contributes at most one new stratum layer.  Conversely, if
a pole of order~$j$ vanishes, the corresponding stratum
collapses and the depth decreases.  To achieve
$r_{\max} = r$ with $r$ finite, one needs non-vanishing
poles at orders $1, \ldots, r$ and vanishing at all orders
$> r$.  For $r > 4$, this means the OPE must have a sharp
cutoff at order~$r$, a condition not satisfied by any
algebra in the standard landscape, where either the pole
order is $\leq 4$ (classes $\mathbf{G}$, $\mathbf{L}$,
and $\mathbf{C}$) or grows without bound (class~$\mathbf{M}$).
\end{proof}

\begin{remark}[Analogy with Coxeter numbers]
% label removed: rmk:gap-coxeter
The gap question is analogous to asking whether there exist
Weyl groups with Coxeter number between~$5$ and~$\infty$
that are not in the standard ADE classification, a
representation-theoretic rigidity question.  Just as the
Coxeter number is constrained by the root system, the
stratification depth is constrained by the FM operadic
structure and the pole data of the OPE.
\end{remark}

\begin{openproblem}[Non-principal complexity]
% label removed: open:thqg-non-principal
\index{non-principal W-algebra!complexity}
For non-principal W-algebras $\mathcal{W}_k(\fg, f)$:
what is the shadow depth class?  The hook-type corridor
in type~A is proved
(Volume~I); the general case
is open.  Expectation: all non-abelian non-principal
W-algebras are class~$\mathbf{M}$ (the DS reduction
introduces composite fields at all arities).
\end{openproblem}

\begin{openproblem}[Categorical complexity]
% label removed: open:thqg-categorical
\index{categorical shadow depth}
Is there a categorical lift of the shadow depth that
detects $r_{\max}$ at the level of derived categories
without passing through explicit OPE computations?
Candidate: the $t$-structure filtration depth on the
derived category of modules over $\cA$, measured by
the Ext-vanishing pattern.
\end{openproblem}

% ----------------------------------------------------------------------
%  The t-structure interpretation (Chriss--Ginzburg attack)
% ----------------------------------------------------------------------

\subsection{The \texorpdfstring{$t$}{t}-structure interpretation}
% label removed: subsec:t-structure-interpretation
\index{t-structure@$t$-structure!amplitude}
\index{shadow depth!$t$-structure detection}

\providecommand{\Steinberg}{\mathfrak{S}}

\begin{conjecture}[Shadow depth as $t$-structure amplitude]
\label{prop:shadow-t-amplitude}
\ClaimStatusConjectured
Let $\cA$ be a modular Koszul chiral algebra with
Steinberg object $\Steinberg_b$.  Then the shadow depth
$r_{\max}$ equals the \emph{amplitude} of the natural
$t$-structure on $D^b(\mathrm{Coh}(\Steinberg_b))$: the
length of the longest non-trivial filtration in the heart.

In classical geometric representation theory, the
perverse $t$-structure on $D^b(\mathrm{Perv}(G/B))$ has
amplitude equal to the length $\ell(w_0)$ of the longest
element of the Weyl group~$W$.  The chiral analogue reads:
\[
  \mathrm{amp}\bigl(D^b(\mathrm{Coh}(\Steinberg_b))\bigr)
  \;=\;
  \max\bigl\{\, k \mid H^k(B(\cA)) \neq 0 \,\bigr\}
  \;=\;
  r_{\max}.
\]
That is, the $t$-structure amplitude on $D^b(\Steinberg_b)$
coincides with the maximal bar degree carrying non-vanishing
bar cohomology.
\end{conjecture}

\begin{remark}[Intrinsic detection of complexity class]
% label removed: rmk:t-structure-intrinsic
Conjecture~\ref{prop:shadow-t-amplitude} provides a
categorical (intrinsic) detection of $r_{\max}$ without
explicit OPE computation.  To determine whether a chiral
algebra belongs to class $\mathbf{G}$, $\mathbf{L}$,
$\mathbf{C}$, or $\mathbf{M}$, one computes the amplitude of the natural
$t$-structure on the derived category of the Steinberg
object, a purely homological-algebraic invariant.
This answers Open Problem~\ref{open:thqg-categorical}
affirmatively, modulo the conjecture above.
\end{remark}

\begin{remark}[Simple objects and stratification]
% label removed: rmk:t-structure-simples
For perverse sheaves on the classical Steinberg variety,
the $t$-structure amplitude equals the number of simple
perverse sheaves, which in turn equals $|W|$.  For the
chiral Steinberg $\Steinberg_b$, the r\^ole of simple
objects is played by the irreducible components of the
orbit stratification; their number is the shadow depth
$r_{\max}$.  The passage from $|W|$ to $r_{\max}$ is the
chiral shadow of the passage from the Weyl group to the
OPE pole pattern.
\end{remark}
